\documentclass{tufte-handout}

\title{MATH 23500 Lecture Notes}

\author{Marcus Kuo}

\date{\today} % without \date command, current date is supplied

%\geometry{showframe} % display margins for debugging page layout

\usepackage{lastpage} % Required to determine the last page number for the footer
\usepackage{enumerate}
\usepackage{amssymb}
\usepackage{verbatim}
\usepackage{mathrsfs}
\usepackage[most]{tcolorbox} % Required for boxes that split across pages
\usepackage{booktabs} % Required for better horizontal rules in tables
\usepackage{listings} % Required for insertion of code
\usepackage{etoolbox} % Required for if statements
\usepackage{mathrsfs}
\usepackage{amssymb}
\usepackage{capt-of}
\usepackage{soul}
\usepackage{pdfpages}
\usepackage{xcolor}
\usepackage{hyperref}
\usepackage{graphicx} % allow embedded images
\setkeys{Gin}{width=\linewidth,totalheight=\textheight,keepaspectratio}
\graphicspath{{}} % set of paths to search for images
\usepackage{amsmath}  % extended mathematics
\usepackage{booktabs} % book-quality tables
\usepackage{units}    % non-stacked fractions and better unit spacing
\usepackage{multicol} % multiple column layout facilities
\usepackage{lipsum}   % filler text
\usepackage{fancyvrb} % extended verbatim environments
\fvset{fontsize=\normalsize}% default font size for fancy-verbatim environments
\usepackage{amsthm}
\usepackage{dsfont}
\usepackage{mathtools}
\DeclarePairedDelimiter{\ceil}{\lceil}{\rceil}

% Standardize command font styles and environments
\newcommand{\doccmd}[1]{\texttt{\textbackslash#1}}% command name -- adds backslash automatically
\newcommand{\docopt}[1]{\ensuremath{\langle}\textrm{\textit{#1}}\ensuremath{\rangle}}% optional command argument
\newcommand{\docarg}[1]{\textrm{\textit{#1}}}% (required) command argument
\newcommand{\docenv}[1]{\textsf{#1}}% environment name
\newcommand{\docpkg}[1]{\texttt{#1}}% package name
\newcommand{\doccls}[1]{\texttt{#1}}% document class name
\newcommand{\docclsopt}[1]{\texttt{#1}}% document class option name

\newcommand*\N{\mathbb{N}}
\newcommand*\R{\mathbb{R}}
\newcommand*\Q{\mathbb{Q}}
\newcommand*\Z{\mathbb{Z}}
\newcommand*\E{\mathbb{E}}
\let\oldP\P
\renewcommand*\P{\mathbb{P}}
\let\O\varnothing
\let\oldemptyset\emptyset
\let\emptyset\varnothing

\DeclareMathOperator*{\argmax}{arg\,max} % thin space, limits underneath in displays
\DeclareMathOperator*{\argmin}{arg\,min} % thin space, limits underneath in displays

\newenvironment{docspec}{\begin{quote}\noindent}{\end{quote}}% command specification environment

\newtheorem{thm}{Theorem}[subsection]
\newtheorem{lemma}[thm]{Lemma}
\newtheorem{cor}[thm]{Corollary}
\newtheorem{rmk}[thm]{Remark}
\newtheorem{conj}[thm]{Conjecture}
\newtheorem{ex}[thm]{Example}
\newtheorem{aside}[thm]{Aside}
\newtheorem{notation}[thm]{Notation}
\newtheorem{prop}[thm]{Proposition}

\theoremstyle{definition}
\newtheorem{defn}[thm]{Definition}

\setcounter{secnumdepth}{2}
\setcounter{tocdepth}{1}

\relpenalty=10000
\binoppenalty=10000

\begin{document}

\maketitle
\tableofcontents

\newpage

\section{Markov Processes, Introduction}
\subsection{Preliminaries}

\begin{defn}[Stochastic process]
	A stochastic process is a collection of random variables $\{X_t; t \in T\}$ where each $x_t$ takes values in $S$. We call $S$ the state space.

	Alternatively, we have a random function $X: T \to S$. $T$ can be discrete ($\N \cup \{0\}$, often notated with $n$ instead of $t$) or continuous, as in $(0,\infty)$. We assume $S$ is countable, except for when it is a continuous space and $\R^d$, where $d \ge 1$.\footnote{For the time being, $T$ and $S$ will be discrete.}
\end{defn}

\begin{ex}
	$\{X_n\}_{n \ge 0}$ is a sequence of independent random variables. $\{Y_n\}_{n \ge 0}$ is a sequence of independent and identically distributed random variables in $\R$. We can have $Z_0 = 0$ and $Z_n = \sum\limits_{j=1}^{n} Y_j$ for $n\ge 1$.
\end{ex}

\begin{defn}[Distribution]
	If $Y$ is a random variable in $S$ and $S$ is countable, the distribution of $Y$ is $y \mapsto \P\left[ Y= y \right]$.

How do we describe the distribution of $\{X_n\}_{n\ge 0}$ as $\P\left[ X_0 = s_0, \ldots, X_n = s_n \right]$ for every $n \in N$ and $s_n \in S$? We cannot say much about any given stochastic process and must impose some structure. 

\end{defn}

\begin{defn}[Conditional probability]
	Let $E,F$ be events such that $\P[F] > 0$. The conditional probability of $E$ given $F$ is \[
		\P[E \mid F] = \frac{P[E \cap F]}{P[F]}
	.\] 
\end{defn}

By definition, for each $n$ and $s_0,\ldots.,s_n \in S$,
\begin{align*}
	\P[X_0=s_0, \ldots, X_n = s_n] &= P[X_n = s_n \mid X_0 = s_0, \ldots, X_{n-1}] = s_{n-1}] \cdot P[X_0 = s_0,\ldots,X_{n-1} = s_{n-1}] \\
	\P[X_n &= s_n \mid \ldots] \cdot P[X_{n-1} = s_{n-1} \ mid \ldots] \cdot \ldots
.\end{align*}

Iterating, the distribution of $\{X_n\}$ is determined by the distribution of $X_0$ and $\P\left[ X_n = sn \mid X_0 = s_0,\ldots,X_{n-1} = s_{n-1} \right]$ for every $n$ and $s_0,\ldots,s_n$.

\begin{defn}[Markov process]
	$\{X_n\}$ is a Markov process (Markov chain) if for every $n$ and $s_0,\ldots,s_n \in S$, \[
		\P\left[ X_n = s_n \mid X_0 = s_0, \ldots, X_{n-1} = s_{n-1} \right] = P\left[ X_n = s_n \mid X_{n-1} = s_{n-1} \right] 
	.\] At time $n$, the only necessary information is the state at time $n-1$ and nothing before.
\end{defn}

\begin{defn}[Time homogeneous]
	$\{X_n\}$ is time homogeneous if for every $n$ and $x,y \in S$, we have \[
		\P\left[ X_n = y \mid X_{n-1} = x \right] = P\left[ X_1 = y \mid X_0 = x \right] 
	.\] In a non-time homogenous Markov process, the conditional probabilities can depend on $n$. 
\end{defn}

All of the Markov process we consider will be time homogeneous.

To specify a (time homogeneous) Markov process, we only need to describe the distribution of $X_0$ and the transition probabilities $\P\left[ x,y \right] := P\left[ X_1 = y \mid X_0 = x \right].$

\begin{ex}
	Let $Y_j$ be iid in $\Z$. Let $X_0 = 0, X_n = \sum\limits_{j=1}^{n} X_j$ for $n \ge 1$. This is a Markov process: $X_n = X_{n-1} + Y_n$. The transition probabilities are given by 
	\begin{align*}
		\P[x,y]  &= \P\left[ X_n \mid y \mid X_{n-1} = x \right]  \\
		&= \P\left[ Y_n = y - x \mid X_{n-1} = x \right] \\
		&= \P\left[ Y_1 = y-x \right] 
	.\end{align*}
\end{ex}

\begin{ex}
	Let $S = \{0,1\}$. We know $\P[0,0] + P[0,1] = 1$ and $P[1,0] + P[1,1] = 1.$ The distribution of the Markov chain is specified by $p = P[0,1]$ and $q = P[1,0]$.
\end{ex}

\begin{defn}[Graph]
	A graph $G$ is a collection of vertices $V(G)$ and edges $E(G)$ joining pairs of vertices. We take the edges to be undirected and allow infinite graphs, with the condition that each vertex is incident to finitely many edges.
\end{defn}

\begin{defn}[Random walk]
	Let $\text{deg}(x)$ denote the number of neighbors of the vertex $x$. A random walk on a graph $G$ iis the Markov chain with \[
		\P[x,y] = \begin{cases}
			\frac{1}{\text{deg}(x)}, &\quad x,y \text{ joined by an edge}, \\
			0 &\quad \text{otherwise}
		\end{cases}
	.\] 
\end{defn}

\begin{ex}
	Let $G = \Z$ with edges joining $x$ and $x+1$ for all $x \in \Z$.
	Let $\{Y_j\}_{j \ge 1}$ to be iid with $\P\left[ Y_j = 1 \right] = P\left[ Y_j = -1 \right] = \frac{1}{2}$. Then $X_n = \sum\limits_{j=1}^{n} Y_j$ is a random walk on the integers.
\end{ex}

\begin{ex}
	Let $S=\Z$ and take $X_0=X_1=X_2 =0$, then \[
		X_n = \begin{cases}
			X_{n-1} + 1 &\quad \text{probability } \frac{1}{3} \\
			X_{n-1} - 1 &\quad \text{probability } \frac{1}{3} \\
			X_{n-3} &\quad \text{probability } \frac{1}{3}
		\end{cases}
	.\] 
	This is not a Markov process because $X_n$ depends on $X_{n-3}$ and not just $X_{n-1}$.
\end{ex}

\begin{defn}[$n$-step transition probability]
	The $n$-step transition probabilities are \[
		P^n(x,y) = \P\left[ X_n = y \mid X_0 = x \right] 
	\] for all $x,y \in S$.
\end{defn}

\begin{prop}
For all $n,m\in\mathbb{N}$ and $x,y\in S$,
\[
P^{n+m}(x,y)=\sum_{z\in S} P^{n}(x,z)\,P^{m}(z,y).
\]
\end{prop}

\begin{proof}
Write $\mathbb{P}_x(\,\cdot\,):=\mathbb{P}(\,\cdot\,\mid X_0=x)$, so that
$P^{k}(a,b)=\mathbb{P}_a(X_k=b)$. Then
\[
P^{n}(x,z)\,P^{m}(z,y)
= \mathbb{P}_x(X_n=z)\,\mathbb{P}_z(X_m=y)
= \mathbb{P}_x(X_n=z)\,\mathbb{P}(X_{n+m}=y\mid X_n=z),
\]
where the last equality uses time homogeneity. By the Markov property,
\[
\mathbb{P}(X_{n+m}=y\mid X_n=z,X_0=x)=\mathbb{P}(X_{n+m}=y\mid X_n=z).
\]
Hence
\[
\mathbb{P}_x(X_n=z)\,\mathbb{P}(X_{n+m}=y\mid X_n=z)
= \mathbb{P}_x\big(X_n=z,\,X_{n+m}=y\big).
\]
Summing over $z\in S$ gives
\[
\sum_{z\in S} P^{n}(x,z)\,P^{m}(z,y)
= \sum_{z\in S} \mathbb{P}_x\big(X_n=z,\,X_{n+m}=y\big)
= \mathbb{P}_x(X_{n+m}=y)
= P^{n+m}(x,y).
\]
\end{proof}
Assume we have a finite state space --- $S$ is finite. WLOG, label the states $s_1,\ldots,s_n$.

\begin{defn}[Transition matrix]
	Define the transition matrix to be the $N \prod N$ matrix $P$ such that $P_{i,j} = P(i,j)$. The $i$th row gives the probabilities for when the state is $s_i$, and the $j$th column gives the probability for moving to the $j$th state. The probabilities of each \emph{row} sum to $1$. We write \[
		P = \begin{pmatrix}
			P(1,1) & \ldots & P(1,N) \\
			P(2,1) & \ddots & \vdots \\
			\vdots & & \\
			P(N,1) & \ldots & P(N,N)
		\end{pmatrix} 
	.\] 

	Put $\pi_j = \P\left[ X_0 = j \right]$, so $\pi := \left( \pi_1,\ldots,\pi_N \right)$ is a row vector encoding the starting probabilities.
\end{defn}

\begin{prop}
	For every $i = 1,\ldots,N$, the $i$th entry of $\pi P$\marginnote{The vector is on the left side of the matrix.} is $\P\left[ X_1 = i \right]$
\end{prop}

\begin{proof}
	Note that 
	\begin{align*}
		\left( \pi P \right)_i &= \sum\limits_{j=1}^{N} \pi_j P(j,i) \\
							   &= \sum\limits_{j=1}^{N} \P[X_0=j] P[X_1 = i \mid X_0 = j] \\
							   &= \sum\limits_{j=1}^{N} \P\left[ X_1 = i, X_0 = j \right]  \\
							   &= \P\left[ X_1= i \right] 
	.\end{align*}
\end{proof}

\begin{prop}
	For every $n \in \N$, $(P^n)_{i,j} = P^n(i,j)$.\marginnote{\texthl{how much computation with linalg are we expected to do on the tests}}
\end{prop}

\begin{proof}
	Induct on $n$.

	Base case ($n=1$): True by definition.
	
	Now assume $(P^{n-1})_{i,j} = P^{n-1}(i,j)$. Then
	\begin{align*}
		\left( P^n \right)_{i,j} &=  \left( P^{n-1}P \right)_{i,j} \\
		&= \sum\limits_{k=1}^{N} \left( P^{n-1}_{i,k} \right) P_{k,j} \\
		&= \sum\limits_{k=1}^{N} P^{n-1}(i,k) P(k,j) \\
		&= P^n(i,j)
	\end{align*} by the previous proposition.
\end{proof}

\begin{ex}
	Let $S = \{0,1\}$, $P(0,1) = \frac{1}{3}$, and $P(1,0) = \frac{1}{2}$. Then \[
		P = \begin{pmatrix}
		\frac{2}{3} & \frac{1}{3} \\
		\frac{1}{2} & \frac{1}{2}
		\end{pmatrix} 
	.\] Suppose we want to find the conditional distribution of $X_3$ given that $X_0 = 0$.

	Taking $P$ to the third power, we have \[
		P^3 = \begin{pmatrix}
		\frac{65}{108} & \frac{43}{108} \\
		\frac{43}{72} & \frac{29}{72}
		\end{pmatrix} 
	.\] So,
	\begin{align*}
		\P\left[X_3 = 1 \mid X_0 = 0\right]  &= \frac{65}{108} \\
		\P\left[X_3 = 0 \mid X_0 = 0\right] &= \frac{43}{108}
	.\end{align*}
\end{ex}

\subsection{Recurrent and Transience (Finite $S$)}

\begin{defn}[Communicate]
	The states $x,y \in S$ if there exists $m,n \ge 0$ such that $P^m(x,y) > 0$ and $P^n(y,x) > 0$. If $x,y$ communicate, we write $x \leftrightarrow y$.
\end{defn}

\begin{prop}
	Communication is an equivalence relation.
\end{prop}

\begin{proof}
	It is reflexive, since $P^0(x,x) = 1 > 0$. It is symmetric, since we can swap $m,n$ for $n,m$. It is transitive, since we can connect with intermediate steps.

	Assume $x \leftrightarrow y$ and $y \leftrightarrow z$. Choose $m,n,l,k$ such that $P^n(x,y)$, $P^m(y,x), P^l(y,z), P^k(z,y) > 0$.
	Then 
	\begin{align*}
		P^{n+l}(x,z) &= \sum\limits_{w\in S}^{} P^n(x,w) P^l(w,z) \\
					 &\ge P^n(x,y) P^l(y,z)
					 &> 0
	.\end{align*}
	Similarly, $P^{k+m}(z,x) > 0$, so $x \leftrightarrow z$.
\end{proof}

\begin{defn}[Communication class, recurrent and transient]
	$x$ and $y$ belong to the same communication class $C \subset S$ iff $x \leftrightarrow y$. The "if" implies that the communication class must be as large as possible.

	\marginnote{This definition only applies to when $S$ is finite.}A communication class $C \subset S$ is recurrent if $P(x,y) = 0$ for all $x \in C, y \not\in C$. It is transient otherwise, as it is possible to "leave" the set.
\end{defn}

\begin{defn}[Irreducible]
	The Markov chain $\{X_n\}$ is irreducible if there is only one communication class. Irreducible Markov chains are automatically recurrent.\footnote{Again, only for finite state spaces.}
\end{defn}

\begin{ex}
	Consider a graph $G$ and a random walk $\{X_n\}$ on $G$. Two vertices $x,y \in V(G)$ are in the same communication graph iff there exists a path of edges between $x$ and $y$. Communication classes in the vertices of a graph are connected components of the graph.
\end{ex}

\begin{ex}
	Let $S = \{1,2,3,4,5\}$ and $P = \begin{pmatrix}
		\frac{1}{5} & \frac{1}{5} & 0 & 0 & \frac{3}{5} \\
		0 & 0 & 0 & 0 & 1 \\
		0 & 0 & \frac{1}{2} & \frac{1}{2} & 0 \\
		0 & 0 & \frac{1}{4} & \frac{3}{4} & 0 \\
		\frac{1}{2} & \frac{1}{4} & \frac{1}{4} & 0 & 0
	\end{pmatrix} $.
	Drawing a graph, we see that $\{3,4\}$ is a recurrent communication class, $\{1,2,5\}$ is a transient communication class.\marginnote{\texthl{Last class, we said we take graphs to be undirected.}}
\end{ex}

\begin{ex}
	Consider a RW on $\{0,\ldots,N\}$ with absorbing boundaries:
	\begin{align*}
		P(x,x+1) &= P(x,x-1) = \frac{1}{2} \\
		P(0,0) &= P(1,1) = 1
	.\end{align*}

	A (transient) communication class is $\{1,\ldots,N-1\}$. But if we include $0$ or $N$, the set is no longer a communication class, since $P(0,x) = P(N,x) = 0$ for at least one (all) $x \in \{1,\ldots,N-1\}$.
	$\{0\}$ and $\{N\}$ are recurrent communication classes.
\end{ex}

\begin{prop}
	Assume $S$ is finite. Let $C$ be a recurrent communication class. Then for every $x,y \in C$, \[
		\P\left[\exists \text{ infinitely many } n \text{ such that } X_n = y \mid X_0 = x\right] = 1
	.\] 	
\end{prop}

\begin{proof}
	Fix $x,y \in C$. Assume $X_0 = x$.
	Since $C$ is a communication class and so $z \leftrightarrow y$, $\forall z \in C$, $\exists n_z \ge 0$ such that $P^{n_z}(z,y) > 0$. Let $n = \max_{z} \{n_z: z \in C\}$, $q = \max_z \{P^{n_z}(z,y) : z \in C\}$. Since $S$ is finite, $q$ is finite and $z$ is positive.

	For $k \in \N$, let $E_k = \{\exists  j \in [n(k-1)+1, nk) \text{ s.t. } X_j = y\}$. Let $s_0,\ldots,s_{nk} \in C$. Then \marginnote{\texthl{ I don't know if the last line is copied correctly.}}
	\begin{align*}
		\P\left[E_{k+1} \mid X_0 = s_0, \ldots, X_{nk} = s_{nk}\right] &= \P\left[E_{k+1} \mid X_{nk} = s_{nk}\right]  \\
		&= \P\left[E_1 \mid X_0 = s_{nk}\right]  \quad \text{by time homogeneity} \\
		&\ge P^{n_{s_{nk}}}(s_{nk},y) \ge q
	.\end{align*}
	We have shown a lower bound on $E_{k+1}$ occurring given any events in the past.

	Let $M,N \in \N$ with $M > N$. Then
	\begin{align*}
		\P\left[E_k \text{ does not occur for any } k \in \{N,\ldots,M\}\right] &= \P\left[\bigcap\limits_{k=N}^{M} E_k^{c}	\right]   \\
		&= \P\left[E_M^c \mid \bigcap\limits_{k=N}^{M-1} E_k^c\right] \P\left[\bigcap\limits_{k=N}^{M-1} E_k^c\right]  \\	
		&\le (1-q) \P\left[\bigcap\limits_{k-N}^{M-1} E_k^c\right] \\
		&\le (1-q)^{M-N} \to 0 \text{ as } M \to \infty \text{ with } N \text{ fixed}
	.\end{align*}

	So, $\P\left[E_k \text{ does not occur for any }k \ge N\right] = 0$. Hence, \[
		\P\left[\exists  j \ge nN \text{ s.t. } X_j = y\right] = 1
	,\] and \[
	\P\left[\exists  \text{ infinitely many }j \text{ s.t. } X_j = y\right] = 1
	.\] 
\end{proof}

\begin{rmk}
	Taking $N = 0$ and $M = nk$, we showed that $\exists  q \in (0,1)$ and $n \ge 1$ such that $\forall  k \in \N, \forall x,y, \in C$, \[
		\P\left[X \text{ hits $y$ before time $nk$} \mid X) = x\right] \ge 1- (1-q)^{k}
	.\] 
	Given $j\ge 1$, choose $k$ such that $j \in \left[ nk, n(k+1) \right]$. Put $c = \frac{-\ln{(1-q)}}{n}$. Then \[
		\P\left[X \text{ hits $y$ before $j$} \mid X_0 = x\right] \ge 1 - (1-q)^{\frac{j}{n}} = 1-e^{-cj}
	.\] As the length of the Markov chain increases, the chance that it does not hit an element decays (at least) exponentially.
\end{rmk}

\begin{prop}
	Assume $S$ is finite. Let $C$ be a transient communication class. With probability $1$, $\{X_n\}$ eventually leaves $C$ and never returns.
\end{prop}

\begin{proof}
	Similar to the previous proof. See the lecture notes. 
\end{proof}

\begin{rmk}
	There exists $c > 0$ such that for every $x \in C$ and $j \in \N$, \[
		\P\left[\{X_n\} \text{ exists $C$ before time } j \mid X_0 = x\right] \ge 1-e^{-cj}
	.\] 
	The chance of staying in a transient communication class decreases exponentially.
\end{rmk}

\subsection{Stopping time}

\begin{defn}[Stopping time]
	A random time $\tau \in \N \cup \infty$ is a stopping time if for every $n \in N$, the event $\{\tau = n\}$ is determined by $X_0, \ldots,X_n$.

	There are some equivalent definitions. We can use the event $\{\tau \le n\}$, since \[
		\{\tau \le n\} = \bigcup\limits_{j=0}^{n} \{\tau=j\}
	\] and \[
		\{\tau = n\} = \{\tau \le n\} \setminus \{\tau \le n-1\}
	.\] We can also use the event $\{\tau > n\}$, since \[
		\{\tau > n\} = \{\tau \le n\}^c
	.\] 
\end{defn}

\begin{ex}
	A non-random time like $\tau = n$ is a stopping time.

	$\tau = \min \{n \ge 0: X_n = x\}$ is a stopping time because $\{\tau \le n\} = \{\exists  j \le n \text{ s.t. } X_j = x\}$. Similarly, $\tau = \{k\text{th smallest time such that } X_j \in A \subset S\}$ is a stopping time.

	The minimum of two stopping times is also a stopping time.
\end{ex}

\begin{ex}
	Let $N \in N, x \in S$ and $\tau$ be the last time $n \le N$ such that $X_n = x$. $\tau$ is not a stopping time because if we just see $X_1,\ldots,X_n$ for $n \le N-1$, we can't tell whether we visit $x$ between $n$ and $N$.
\end{ex}

\subsection{Strong Markov Property}

Let $\{X_n\}$ be a Markov chain with a countable state space $S$. For each $x_0,\ldots,x_n \in S$ and $y_1,\ldots,y_m \in S$, we have 
\begin{align*}
	&\P\left[X_{n+1} = y_1,\ldots,X_{n+m} = y_m \mid X_0 = x_0,\ldots,X_n = x_n\right] \\
	&= \P\left[X_{n+1} = y_1, \ldots, x_{n+m} = y_m \mid X_n = x_n\right]  \\
	&= \P\left[x_n, y_1\right] \prod_{i=2}^{m} \P\left[y_{i-1}, y_i\right] 
.\end{align*}

\begin{defn}[Stopping time]
	A random time $\tau \in \N_0 \cup \{\infty\}$ is a stopping time if for every $n \in \N_0$, the event $\{\tau = n\}$ is determined by $X_0,\ldots,X_n$.\marginnote{As a general rule, the "first time" something happens is a stopping time, while the "last time" something happens is \emph{not} a stopping time.}
\end{defn}

\begin{thm}
	(Strong Markov Property) Let $\tau$ be a stopping time. Let $n \ge 0, m \ge 1$, $X_0,\ldots,X_n \in S$ such that $\P\left[X_0= x_0,\ldots,X_\tau = x_n\right] > 0$. Let $y_1,\ldots,y_m \in S$. Then \[
		\P\left[X_{\tau+1} = y_1,\ldots,X_{\tau+m} = y_m \mid X_0 = x_0, \ldots, X_\tau = x_n\right] = \P\left[X_1=y_1,\ldots,X_m = y_m \mid X_0=x_n\right] 
	.\] The defining property of a Markov chain applies to when you end at a stopping time instead of a non-random time.
\end{thm}

\begin{proof}
	Note that \begin{align*}
		\{X_0 = x_0,\ldots,X_\tau = X_n\} &= \{\tau = n\} \cap \{X_0 = x_0, \ldots, X_n = x_n\} \\
										  &= \{X_0 = x_0,\ldots,X_n = x_n\}
	.\end{align*}
	Further,
	\begin{align*}
		&\P\left[X_{\tau + 1} = y_1, \ldots, X_{\tau + m } = y_m \mid X_0 = x_0, \ldots, X_\tau = x_n\right] \\
		&= \P\left[X_{n+1} = y_1,\ldots,X_{n+m} = y_m \mid X_0 = x_0, \ldots, X_n = x_n\right]  \quad \text{since we require that } \tau = n \\
		&= \P\left[X_1 = y_1,\ldots,X_m = y_m \mid X_0 = x_n\right] \quad \text{by the ordinary Markov property}
	.\end{align*}
\end{proof}

\begin{ex}
	Let $x \in S$. Let $\tau = \min \{n \ge 0 : X_n = x\}$. Assume $\P\left[\tau < \infty\right] = 1$.\footnote{Unrelatedly, this implies that the Markov chain is irreducible.} For any $y_1,\ldots,y_m \in S$ and $x_1,\ldots,x_n \in S$ such that $\P\left[X_0 = x_0, \ldots,X_\tau = x_n\right] > 0$. Note that by our definition of $\tau$, this implies that $x_n= x$. (At time $\tau$, $\{X_n\}$ is at $x$.)

	Applying the Strong Markov Property, 
	\begin{align*}
		&\P\left[X_{\tau + 1} = y_1,\ldots,X_{\tau + m} = y_m\mid X_0 = x_0,\ldots,X_\tau = x_n\right] \\
		&= \P\left[X_1 = y_1, \ldots,X_m  = y_m \mid X_0 = x\right] 
	.\end{align*}

	The same is true for any stopping time $X_\tau = x$. For example, the $k$th visit to $x$ or the first visit to $x$ after visiting $y$ could definite $\tau$.
\end{ex}

\begin{prop}
	Suppose $X_0 = x \in S$. Assume $\P\left[\{X_n\} \text{ visits } x \text{ infinitely often}\right]  = 1$. Let $\tau_k = k\text{th time } n \text{ such that } X_n = x$. By convention, set $\tau_0 = 0$.

	Then the increments $\{\left( X_{\tau_k},\ldots,X_{\tau_{k+1}} \right)\}_{k \in \N_0}$ are independent and identically distributed. \marginnote{$\left( X_{\tau_k},\ldots,X_{\tau_{k+1}} \right) \in \bigcup\limits_{j=1}^{\infty} S^j$.}
\end{prop}

\begin{proof}
	Each $\tau_k$ is a stopping time, and $X_{\tau_k} = x$ for every $k$. This implies that the condition distribution of $\{X_{\tau_k} + j\}_{j \ge 0}$ given everything before time $\tau_k$ is the same as $\{X_j\}_{j\ge 0}$ since we started the Markov chain at state $x$. 

	Also observe that $\tau_{k+1} - \tau_k$ is the first time for $j \ge 1$ that $x_{\tau_k + j} = x$. The condition distribution of $\left(X_{\tau_k},\ldots,X_{\tau_{k+1}}\right)$ given everything before time $\tau_k$ is the same distribution of $\left( X_0,\ldots,X_{\tau_1} \right)$.

	This implies that the increments are i.i.d.
\end{proof}

\begin{defn}[Geometric distribution]
	A random variable $M \in \N$ has the geometric distribution with success probability $p \in (0,1)$ if \[
		\P\left[M = m \right]  = p(1-p)^{m-1}
	\] for all $m \ge 1$.
\end{defn}

\begin{prop}
	Suppose $X_0 = x \in S$ and $\P\left[\{X_n\} \text{ visits } x \text{ infinitely often}\right] = 1$. Let $y \in S$ such that $x \leftrightarrow y$.

	Let $M$ be the number of times $\{X_n\}$ visits $x$ before the first time $\{X_n\}$ visits $y$.

	Then $M$ has a geometric distribution.
\end{prop}

\begin{proof}
	Let $\tau_k$ be the $k$th time $n$ such that $X_n = x$. By the previous proposition, $\{\left( X_{\tau_k}, \ldots, X_{\tau_{k+1}} \right)\}_{k \ge 0}$ are i.i.d. Note that $M$ is the smallest $k$ such that $\left( X_{\tau_k},\ldots,X_{\tau_{k+1}} \right)$ visits $y$. $M$ is the first success in a sequence of i.i.d. trials, so it has a geometric distribution.\marginnote{\texthl{what is sufficient to prove something has a geometric distribution?}}
\end{proof}

\subsection{Periodicity}
For this section, let $\{X_n\}$ be a Markov chain with a countable state space $S$.

\begin{defn}[Period]
	For $x \in S$, the period of $x$ is \[
		d(x) = \text{g.c.d.}(J_x)
	,\] where \[
		J_x = \{n \ge 1: P^{n(x,x) > 0}\}
	.\] $J_x$ is the set of times at which we could return to $x$ if we start at $x$, and $d(x)$ is the g.c.d. of this.
\end{defn}

\begin{rmk}
	If $p(x,x) > 0$, then $d(x) = 1$.
\end{rmk}

\begin{defn}[Bipartite]
	A graph $G$ is bipartite if \[
		V(G) = V_1 \cup  V_2
	,\] where $V_1 \cap V_2 = \O$, $V_1,V_2 \ne \O$, and every edge goes from $V_1$ to $V_2$.
\end{defn}

\begin{ex}
	A random walk on $\Z$ is bipartite with $V_1$ the odd integers and $V_2$ the even integers.
\end{ex}

If $G$ is connected and bipartite, then $d(x) = 2$ for a random walk on $G$ for every $x \in V(G)$. If $G$ is connected and not bipartite, then $d(x) = 1$ for every $x \in V(G)$.

\begin{lemma}
	If $n,m \in J_x$, then $n_m \in J_x$.
\end{lemma}

\begin{proof}
	Note that
	\[
		p^{n+m}(x,x) \ge p^{n}(x,x)p^{m}(x,x) > 0
	.\]
\end{proof}

\begin{prop}
	The set $J_x$ contains $kd(x)$ for all sufficiently large $k$.\marginnote{\texthl{intuition for what "sufficiently" means?}}
\end{prop}

\begin{proof}
	The proof requires some number theory and relies on the fact that $J_x$ is closed under addition along with Bezout's identity.
\end{proof}

\begin{prop}
	If $x \leftrightarrow y$, then $d(x) = d(y)$. Consequently, the period of a state only depends on the state's communication class.
\end{prop}

\begin{proof}
	Since $x \leftrightarrow y$, we can choose $n,m$ such that $p^n(x,y) > 0$ and $p^m(y,x) > 0$.

	Then $p^{n+m}(x,x) \ge p^{n}(x,y)p^m(y,x) > 0$.
	Similarly, $p^{m+n}(y,y) > 0$. 

	Since $n_m \in J_x \cap J_y$, both $d(x)$ and $d\left(y \right)$ divide $n+m$.

	Assume for contradiction that $d(x) < d(y)$. Then there exists $k \in J_x$ that is not divisible by $d(y)$. Observe that
	\begin{align*}
		p^{n+m+k}(y,y) &\ge p^m(y,x)p^k(x,x)p^n(x,y) \\
					   &> 0
	,\end{align*} so $n+m+k \in J_y$. But since $n+m \in J_y$ and $d(y)$ divides $n+m$ and $n+m+k$, it also divides their difference, which is $k$. This is a contradiction to our assumption, so $d(x) = d(y)$ by repeating a symmetric argument for $d(y) < d(x)$.
\end{proof}

\begin{defn}[Aperiodicity]
	$\{X_n\}$ is aperiodic if $d(x) = 1$ for all $x \in S$.
\end{defn}

\begin{rmk}
	If $\{X_n\}$ is irreducible, we only need to check $d(x) = 1$ for a single $x \in X$ to check aperiodicity, since the period is dependent only on communication class.
\end{rmk}

\begin{rmk}
	If $\{X_n\}$ is aperiodic, then $p^n(x,x) > 0$ for all sufficiently large $n$ (depending on $x$).
\end{rmk}

\begin{rmk}
	Consider $\tilde{p}(x,y) = \begin{cases}
		\frac{1}{2} p(x,y), \quad x \ne y \\
		\frac{1}{2} + \frac{1}{2}p(x,x), \quad x=x
	\end{cases}$. This is a lazy Markov chain that is aperiodic. If we are given a periodic Markov chan, we can make it aperiodic by slightly modifying it in this way.
\end{rmk}

\subsection{Stationary Distributions (Finite $S$)}

Let $\{X_n\}$ be a Markov chain on state space $S$.

\begin{defn}[Stationary Distribution]
	Let $\pi: S \to [0,1]$ such that $\sum\limits_{x \in S}^{} \pi_x = 1$. $\pi$ is a stationary distribution (invariant distribution) if for every $y \in S$, \[
		\sum\limits_{x \in S}^{} \pi_x p(x,y) = \pi_y
	.\] 

\end{defn}

\begin{rmk}
We have some equivalent definitions.
	\begin{itemize}
		\item (Equivalent definition 1) Suppose $S = \{1,\ldots,N\}$ and $\pi = (\pi_1,\ldots,\pi_N)$ is a row vector. Then $\pi$ is a stationary distribution iff $\sum\limits_{j=1}^{N} \pi_j = 1$ and $\pi P = \pi$ is a left eigenvector of $P$ with eigenvalue $1$.


		\item (Equivalent definition 2) If $\P\left[X_0 = x\right] = \pi_x$ for every $x \in S$, then $\P\left[X_1 = y\right] = \pi_y$ for every $y \in S$.

		Proof: If $\pi$ is a stationary distribution and $X_0 \sim \pi$, then 
			\begin{align*}
				\P\left[X_1 = y\right]  &= \sum\limits_{x \in S}^{} \pi_x p(x,y) \\
										&= \pi_y \quad \text{iff $\pi$ is a stationary distribution}
			.\end{align*}

		\item (Equivalent definition 3) If $\pi$ is a stationary distribution and $\P\left[X_0 = x\right] = \pi_x$ for every $x \in S$ then $X_n \sim \pi$ for every $n \in \N$.
	\end{itemize}
\end{rmk}

\begin{thm}
	Assume $S$ is finite and $\{X_n\}$ is irreducible and aperiodic. Then there exists a unique stationary distribution $\pi$ for every $x,y \in S$ where \[
		\lim\limits_{n \to \infty} \underbrace{\P\left[X_n = y \mid X_0 = x\right] }_{p^{n}(x,y)} = \pi_y
	.\] 
\end{thm}

Note that $\pi_x$ is the long-run average portion of time such that $X_n = x$. We will give a probabilistic proof of this theorem, but it can also be proven using linear algebra (using the left-eigenvector characterization and the Perron-Frobenius Theorem). Furthermore, this theorem is useful for sampling algorithms, since you can run a Markov chain whose stationary distribution matches a given probability distribution $\pi$. 

\begin{ex}
	Let $S = \{1,2,3\}$. Define \[
		P = \begin{pmatrix}
			\frac{1}{3} & \frac{1}{3} & \frac{1}{3} \\
			\frac{1}{2} & 0 & \frac{1}{2} \\
			\frac{1}{5} & \frac{1}{5} & \frac{3}{5}
		\end{pmatrix} 
	.\] The Markov chain is irreducible and aperiodic. Find the stationary distribution such that $\pi P = \pi$ with $\pi_1 + \pi_2 + \pi_3 = 1$. 

	We can compute that $\pi - \left( \frac{3}{10}, \frac{1}{5}, \frac{1}{2} \right)$.
	In the long run, it will spend about $\frac{3}{10}$s of its time in state 1, $\frac{1}{5}$ of its time in state 2, and $\frac{1}{2}$ of its time in state 3.
\end{ex}

\begin{prop}
	Fix $z \in S$ and assume that $X_0 = z$. Let $T = \min \{n \ge 1 : X_n = z\}$. For $x \in S$, let \[
		\tilde{\pi}_x = \E\left[\text{number of } \{n \in \{0,\ldots,T-1\}\} : X_n = x\right] 
	.\] Put $\pi_x = \frac{\tilde{\pi}_x}{\E\left[T\right] }$. Then $\pi$ is a stationary distribution.
\end{prop}

\begin{proof}
	Observe that
	\begin{align*}
		\sum\limits_{x \in S}^{} \tilde{\pi}_x &= 1
	,\end{align*} so the only remaining part to prove is that \[
	\tilde{\pi}_y = \sum\limits_{x \in S} \tilde{\pi}_x p(x,y)
	.\] 

	Note
	\begin{align*}
		\tilde{\pi}_y &= \E\left[\sum\limits_{n=0}^{T-1} \mathds{1}_{\left( X_n = x \right)}\right]  \\
					  &= \E\left[\sum\limits_{n=0}^{\infty} \mathds{1}_{\left( X_n = s, T > n \right)}\right]  \\
					  &= \sum\limits_{n=0}^{\infty} \E\left[\mathds{1}_{\left( X_n = x, T > n \right)}\right]  \\
					  &= \sum\limits_{n= 0}^{\infty} \P\left[X_n = x, \underbrace{T > n}_{\text{depends only on $X_0,\ldots,X_n$}}\right]
	.\end{align*}
	Furthermore for $y \ne z$, \marginnote{\texthl{why can we add in $T > n$? }}
	\begin{align*}
		\sum\limits_{x \in S}^{} \tilde{\pi}_x p(x,y) &= \sum\limits_{x \in S}^{} \sum\limits_{n=0}^{\infty} \P\left[X_n = x, T >n\right] p(x,y) \\
													  &= \sum\limits_{n=0}^{\infty} \sum\limits_{x \in S}^{} \P\left[X_n = x, T > n\right] \P\left[X_{n+1} = y \mid X_n = x, T > n\right]  \\
													  &= \sum\limits_{n=0}^{\infty} \sum\limits_{x \in S}^{} \P\left[X_n = x, T > n, X_{n+1} = y\right]  \\
													  &= \sum\limits_{n=0}^{\infty} \P\left[T > n, X_{n+1} = y\right] \\
													  &= \sum\limits_{n=0}^{\infty} \P\left[T > n+1, X_{n+1} = y\right]  \\
													  &= \sum\limits_{n=1}^{\infty} \P\left[T > n, X_n = y\right] \\
													  &= \tilde{\pi}_y - \P\left[T > 0, X_n = y\right]  \\
													  &= \tilde{\pi}_y
	.\end{align*}\marginnote{\texthl{isn't the last probability that disappears 1?}}

	For $y = z$, we have that 
	\begin{align*}
		\sum\limits_{x \in S}^{} \tilde{\pi}_x p(x,z) &= \sum\limits_{n=0}^{\infty} \P\left[T > n, X_{n+1} = z\right]  \\
		&= \sum\limits_{n=0}^{\infty} \P\left[T = n+1\right]  \\
		&= 1 = \tilde{\pi}_z
	.\end{align*}
\end{proof}

\begin{prop}
	If $\pi$ is a stationary distribution, then for any $x,y \in S$, \[
		\lim\limits_{n \to \infty} \P\left[X_n =y \mid X_0 = x\right] = \pi_y
	.\] In particular, $\pi$ is unique.
\end{prop}

\begin{proof}
	Consider the Markov chain $(X_n, Y_n)$ on $S \times S$ with transition probabilities \[
		\overline{P}\left( \left( x,y \right), \left( x',y' \right) \right) = \begin{cases}
			P(x,y) P(x',y'), \quad x \ne y \\
			P(x,x'), \quad x =y, x'=y' \\
			0, \quad \text{otherwise}
		\end{cases}
	.\] 

	If $x \ne y$, 
	\begin{align*}
		\P\left[X_1 = x' \mid X_0 = x, Y_0 = y\right] &= \sum\limits_{y' \in S}^{} p(x,x')p(y,y') \\
													  &= p(x,x')
	.\end{align*}

	If $x=y$, \[
		\P\left[X_1 = x' \mid X_0 = x, Y_0 = y\right] = p(x,x')
	.\] In works, $\{X_n\}$ and $\{Y_n\}$ have the same transition probabilities as the original Markov chain.

	Let $\tau = \min \{n \ge 0 : X_n = Y_n\}$. By definition, $X_n = Y_n$ for every $n \ge \tau$. We claim that $\P\left[\tau < \infty \mid X_0= x, Y_0 = y\right] = 1$ for every $x,y \in S$.

	Consider the initial distribution where $X_0 = x \in S$ and $Y_0 \sim \pi$. Since $\pi$ is a stationary distribution, $Y_n \sim \pi$ for every $n \in \N$ and if the claim holds, $X_n = Y_n$ for any large enough $n$; i.e. $\lim\limits_{n \to \infty} \P\left[X_n = Y_n \right] = 1$

	This implies that
	\begin{align*}
		&\lim\limits_{n \to \infty} \P\left[X_n = y \mid X_0 = x\right] - \pi_y \\
		&= \lim\limits_{n \to \infty} \P\left[X_n = y \mid X_0 = x\right]  - \P\left[Y_0 =y\right] = 0
	.\end{align*}
	
	It remains to actually prove the claim. Consider $\{\tilde{X}_n\}$ and $\{\tilde{Y}_n\}$ that are independent copies of the Markov chain with $\tilde{X}_0 = x$ and $\tilde{Y}_0 = y$. It suffices to show that $\P\left[\exists n \text{ s.t. } \tilde{X}_n = \tilde{Y}_n\right]$. Hence it is enough to show that $\left( \tilde{X}_n, \tilde{Y}_n \right)$ is an irreducible Markov chain in $S \times S$.

	The aperiodicity assumption implies that there exists $k_0 \in \N$ such that for any $k > k_0$ and $x \in S$, $p^k(x,x) > 0$. The irreducibility assumption implies that for any $x,x',y,y' \in S$, there exists $n \in \N$ such that $p^n(x,x') >0$ and $m \in \N$ such that $m \ge n + k_0$\footnote{Since an irreducible Markov chain will visit each state infinitely many times} and $p^m(y,y') > 0$.

	Since $\tilde{X}_m$ and $\tilde{Y}_m$ are independent by definition, \begin{align*}
		\P\left[\left( \tilde{X}_m, \tilde{Y}_m \right) = \left( x',y' \right) \mid \left( \tilde{X}_0, \tilde{Y}_0 \right) = (x,y)\right] &= p^m(x,x')p^m(y,y') \\
														&\ge p^{\overbrace{m-n}^{\ge k_0}}(x,x)p^{n}(x,x')p^{m}(y,y') \\
														&> 0
	.\end{align*}

	Hence $\left( \tilde{X}_n, \tilde{Y}_n \right)$ is irreducible, so with probability 1, there exists $n$ such that $\left( \tilde{X}_n, \tilde{Y}_n \right) = (x,x)$.
\end{proof}

Let $\{X_n\}$ be a Markov chain on a finite state space $S$ with stationary distribution given by $\pi$ satisfying $\sum\limits_{x \in S}^{} \pi_x = 1$, $\pi_y = \sum\limits_{x \in S}^{} \pi_x p(x,y)$. If $\{X_n\}$ is irreducible and aperiodic, then there is a unique stationary distribution that is the limit of iterating the transition probabilities: $\pi_y = \lim\limits_{n \to \infty} p^n (x,y)$.

\begin{prop}
	Let $\pi$ be the unique stationary distribution of $\{X_n\}$. For $x \in S$, put $T_x = \min \{n \ge : X_n = x\}$. Then \[
		\pi_x = \frac{1}{\E\left[T_x \mid X_0 = x\right]}
	.\] 
\end{prop}
\begin{proof}
	Assume $X_0 = x$. Then \[
		\pi_y = \frac{1}{\E\left[T_x\right]} \E\left[\# \{n \in \{0,\ldots,T_x - 1\} : X_n = y\}\right] 
	.\] If $y=x$, this gives $\pi_y = \frac{1}{\E\left[T_x\right]}$.
\end{proof}

\begin{ex}
	Let $G = (V,E)$ be a finite, not bipartite graph, \texthl{hence (?)} connected.\marginnote{The degree of $x$ is the number of edges connected to $x$.} The stationary distribution for a random walk on $G$ is given by \[
		\pi_x = \frac{deg(x)}{2 \left| E \right|}
	.\] 

	To see why, note for $x \in V$,
	\begin{align*}
		\sum\limits_{x \in V}^{} \pi_x &= \frac{1}{2 \left| E \right|} \sum\limits_{x \in V}^{} deg(x) \\
									   &= 1
	,\end{align*} and for $y \in V$,
	\begin{align*}
		\sum\limits_{x \in V}^{} \pi(x,y) &= \sum\limits_{x \in V, x \sim y}^{} \left( \frac{deg(x)}{2 \left| E \right|} \right) \cdot \frac{1}{deg(x)} \\
										  &= \frac{deg(y)}{2 \left| E \right|} \\
										  &= \pi_y
	.\end{align*}
\end{ex}

\subsection{Recurrent and Transience (Countable $S$)}

Now suppose $S$ is countably infinite.

\begin{defn}[Irreducible]
	$\{X_n\}$ is irreducible if for every $x,y \in S$, there exists $n \in \N$ such that $p^n(x,y) > 0$.
\end{defn} \marginnote{\texthl{ is there any difference between this definition and the communication class definition we used for the finite state space case? This is for a state, the other is for a communication class.}}

\begin{defn}[Recurrent]
	A state $x\in S$ is recurrent if \[
		\P\left[\text{there exist infinitely many n such that } X_n  = x \mid X_0 = x\right]  = 1
	.\] 
	$x$ is transient otherwise.
\end{defn}

While on a finite state space, irreducibility implied recurrence, this is not true for an infinite state space.

\begin{prop}
	If $\{X_n\}$ is irreducible then either every $x \in S$ is recurrent or every $x \in S$ is transient.

	Hence, it makes sense to say that a Markov chain is recurrent or transient.
\end{prop}

\begin{proof}
	Assume there exists a recurrent state $x$. Let $\tau_0 = 0$ and $\tau_k$ be the $k$th smallest $n$ such that $X_n = x$. By our assumption, $\P\left[\tau_k < \infty\right] = 1$ for every $k$. By the Strong Markov Property, $\left( X_{\tau_k},\ldots,X_{\tau_{k+1}} \right) \in \bigcup\limits_{n \in N}^{} S^{n}$ are i.i.d.

	Let $y \in S$. Since the Markov chain is irreducible, there exists $n \in \N$ such that $p^n(x,y) > 0$. So, there exists $k$ such that $\P\left[y \in \{X_{\tau_k}, \ldots,X_{\tau_{k+1}}\}\right] > 0$. Because these tuples (sets) are i.i.d, the probabilities do not depend on the choice of $k$, so we can put \[
		q := \P\left[y \in \{X_{\tau_k},\ldots, X_{\tau_{k+1}}\}\right] 
	.\] 

	Varying $k$, the events $\{y \in \{X_{\tau_k}, \ldots,X_{\tau_{k+1}}\}\}$ are i.i.d and have probability $q > 0$. So, with probability $1$, infinitely many of these events occur; i.e. \[
		\P\left[\text{there exist infinitely many n such that } X_n = y \mid X_0 = x\right] = 1
	.\]

	Let $\sigma = \min \{n \ge 0: X_n = y\}$. We know that $\P\left[\sigma < \infty\mid X_0 = \times \right] = 1$ and by the Strong Markov Property, $\{X_{\sigma + j}\}_{j \ge 0}$ has the same distribution as $\{X_j\}_{j\ \ge_0}$ started at $y$.

	Hence, the Markov chain at $y$ is visits $y$ infinitely often.

\end{proof}

How can we tell if a Markov chain with a countable state space is recurrent or transient? 

\begin{prop}
	A state $x$ is recurrent iff \[
		\sum\limits_{n=0}^{\infty} p^n(x,x) = \infty
	.\] Furthermore, if $x$ is transient (the sum is finite), then with probability $1$, $\{X_n\}$ visits $x$ only finitely many times.
\end{prop}

\begin{proof}
	Let $R_x = \sum\limits_{n=0}^{\infty} \mathds{1}_{\left( X_n = x \right)}$.

	Note
	\begin{align*}
		\E\left[R_x \mid X_0 = x\right] &= \sum\limits_{n=0}^{\infty} \P\left[X_n = x \mid X_0 = x\right]  \\
										&= \sum\limits_{n=0}^{\infty} p^n(x,x)
	.\end{align*}

	As an aside, if we have $\sum\limits_{n=0}^{\infty} p^n(x,x) < \infty$, then $\E\left[R_x\right] < \infty$, so $\P\left[R_x < \infty\right] = 1$ and the Markov chain is transient. The inverse is not true --- the Cauchy random variable has infinite expectation but is bounded with probability $1$.

	Assume that $x$ is transient. We need to show that $\sum\limits_{n=0}^{\infty} p^n(x,x) < \infty$.
	Let $\tau_0 = 0$ and put $\tau_k$ be the $k$th time that $X_n = x$.
	Since $x$ is transient, there exists $k$ such that $\P\left[\tau_k = \infty\right] > 0$.
	By the Strong Markov Property, for every $k$, \[
		\P\left[\tau_{k+1} - \tau_k = \infty \mid \tau_k < \infty\right] =: q > 0
	.\] 

	Note that $R_x =\min \{k : \tau_{k+1} = \infty\}$. Then $R_x$ has a geometric distribution with success probability $q$, so \[
		\sum\limits_{n=0}^{\infty} p^n(x,x) = \E\left[R_x \mid X_0= x\right] = \frac{1}{q} < \infty
	.\] 

\end{proof}

If $\{X_n\}$ is irreducible and we want to check if it is recurrent or transient, then we just need oto check if $\sum\limits_{n=0}^{\infty} p^n(x,x) = \infty$ for some $x \in S$. \marginnote{\texthl{If the Markov chain is reducible, then it is automatically transient.}}

\begin{ex}
	Take $S = \N$ and $p(x,0) = \frac{1}{x+2}, p(x,x+1) = 1 - \frac{1}{x+2}$. 
	This Markov chain is irreducible since from every state, we have a positive probability of going to 0, and from 0 we have positive probability of going to every other state.

	To determine if it is recurrent or transient, assume $X_0 = 0$. Then $X_n \le n$.

	This implies that $P(X_n = 0) \mid X_{0} = 0 \ge \frac{1}{n+1}$. Hence, $p^n(0,0) \ge \frac{1}{n+1}$ and since $\sum\limits_{n=0}^{\infty} p^n(0,0) = \infty$, the Markov chain is recurrent.

	The property of recurrence was very dependent on the transition probability; if we had $p(x,0) = \frac{1}{\left( x+2 \right)^2}$, then the Markov chain would be transient.
\end{ex}

\begin{ex}
	Consider the same Markov chain as the previous example. Let $\tau = \min \{n \ge 1 : X_n = 0\}$.
	To show that the Markov chain is recurrent, it suffices to show that $P\left[ \tau < \infty \right] = 1$, so it suffices to show that $\lim\limits_{N \to \infty} \P\left[\tau > N\right] = 0$.

	We have that $\tau > N$ if the first $N$ steps are upward. Hence,
	\begin{align*}
		\P\left[\tau > N\right] &= \prod_{n=0}^{N} p(n,n+1) \\
								&= \prod\limits_{n=0}^{N} (1-\frac{1}{n+2}) \\
								\ln{(\P\left[\tau > n\right] )} &= \sum\limits_{n=0}^{N} \ln{(1- \frac{1}{n+2})} \\
																&= - \sum\limits_{n=0}^{N} \frac{1}{n+2} +  O\left( \frac{1}{\left( n+2 \right)^2} \right) \\
																&\to -\infty \quad \text{as } N \to \infty \\
								\implies &\P\left[\tau > N\right] \to e^{-\infty}  =0
	.\end{align*}
\end{ex}

\subsection{Biased Random Walk}

\begin{defn}[Biased random walk]
	The biased random walk on $\Z$ is given by fixing $p \in (0,1)$ and putting $p(x,x+1) = p$ and $p(x,x-1) = 1=p$. 
\end{defn}

\begin{prop}
	Let $N \ge 1$. For $x \in \{0,\ldots,N\}$, \[
		\P\left[\{X_n\} \text{ hits } N \text{ before } 0 \mid X_0 = x\right] = \begin{cases}
		\frac{\left( \frac{1-p}{p} \right)^{x} - 1}{\left( \frac{1-p}{p} \right)^{N} - 1}, \quad p \ne \frac{1}{2} \\
		\frac{x}{N}, \quad p = \frac{1}{2}
		\end{cases}
	.\] 
\end{prop}

\begin{proof}
	Define $\alpha(x) = \P\left[\{X_n\} \text{ hits } N \text{ before } 0 \mid X_0 = x\right]$. Note $\alpha(0) = 0$ and $\alpha(N) = 1$.

	Furthermore, $\alpha(x) = p \alpha(x+1) + (1-p) \alpha(x-1)$ for every $x \in \{ 1, \ldots,N-1\}$. We have $N+1$ equations and $N+1$ unknowns.

	If $p=\frac{1}{2}$, we have \[
		\alpha(x) = \frac{\alpha(x+1) + \alpha(x-1)}{2}
	\] for every $x \in \{1,\ldots,N-1\}$, so $\alpha$ is linear. Given our boundary conditions, we have that $\alpha(x) = \frac{x}{N}$.

	If $p \ne \frac{1}{2}$, then we conjecture $\alpha(x) = b^{x}$, so
	\begin{align*}
		b^{x} &= pb^{x+1} + (1-p)b^{x-1} \\
		b &= pb^2 + 1-p \\
		b &= 1 \text{ or } b = \frac{1-p}{p}
	.\end{align*}

	The general solution is then a linear combination of these two specific solutions, so \[
		\alpha(x) = c_1 + c_2 \left(\frac{1-p}{p}\right)^{x}
	.\] Plugging in our boundary conditions, we get that $c_1 = -c_2 = \left( 1- \left( \frac{1-p}{p} \right)^{N} \right)^{-1}$.
\end{proof}

\begin{prop}
	For $x \ge 1$, \[
		\P\left[\{X_n\} \text{ hits } 0 \mid X_0 = x\right] = \begin{cases}
			1, \quad p \le \frac{1}{2} \\
			\left( \frac{1-p}{p} \right)^{x}, \quad p > \frac{1}{2}
		\end{cases}
	.\] 	
\end{prop}

\begin{proof}
	Send $N \to \infty$ in the previous proposition.
\end{proof}

\begin{cor}
	A biased random walk in recurrent if $p = \frac{1}{2}$ and transient if $p \ne \frac{1}{2}$.
\end{cor}

\begin{proof}
	If $p > \frac{1}{2}$, there is a positive chance to never hit $0$, so it is transient. If $p < \frac{1}{2}$, then $\{-X_N\}$ is a biased random walk with $1-p$ in place of $p$. Thus, $\{-X_n\}$ is transient, so $\{X_n\}$ is transient.

	If $p=\frac{1}{2}$, then $\{X_N\}$ has probability $1$ of hitting 0 regardless of $X_0$. In particular, there is probability $1$ of returning to 0 if the random walk starts at 0, so it is recurrent.
\end{proof}

\subsection{Queuing Model}

Suppose at each time $n \ge 1$,  if there is at least one person in line, then one person is served with probability $q \in (0,1)$. After this event, with probability $p \in (0,1)$, a new person arrives. Let these probabilities be independent from each other and the past. Let $X_n$ be the number of people in the queue at time $n$, so the natural choice is $S = \N$. Note that $p(0,1) = p$ and $p(0,0) = 1-p$. Furthermore, for $x \ge 1$, $p(x,x-1) = q(1-p)$ and $p(x,x+1) = p (1-q)$. Finally, we have that $p(x,x) = qp + (1-q)(1-p)$.

This Markov chain is irreducible. We wonder if the Markov chain is recurrent or transient.

\begin{prop}
	$\{X_n\}$ is recurrent iff $q \ge p$. 
\end{prop}
The following proof technique is useful --- we reduce this Markov chain to a biased random walk.
\begin{proof}
	Let $\tau_k$ be the $k$th time such that $X_n \ne X_{n-1}$. Note that each $\tau_k$ is a stopping time. 
	We have that\marginnote{The first equality is because the Markov chain does not move between $\tau_k and \tau_{k-1}$ and so the transition probability to $x+1$ is the same.}
	\begin{align*}
		\P\left[X_{\tau_k} = x+1 \mid X_{\tau_{k-1}} = x\right] &= \P\left[X_1 = x+1 \mid X_0 = x, X_1 \ne X_0\right]  \\
																&= \frac{p(x,x+1)}{1-p(x,x)} \\
																&= \frac{p(1-q)}{q(1-p) + p(1-q)}
	.\end{align*}

	Since the Markov chain does not move between $\tau_{k-1}$ and $\tau_{k}$, we have that
	\begin{align*}
		\P\left[X_{\tau_k} = x-1 \mid X_{\tau_{k-1}} = x\right] = 1 - \frac{p(1-q)}{q(1-p) + p(1-q)} 
	.\end{align*}
	$\{X_{\tau_k}\}$ is a biased random walk with parameter $\frac{p(1-q)}{q(1-p) + p(1-q)}$ (until it hits 0). We see that $\{X_{\tau_k}\}$ eventually hits 0 with probability $1$ iff \[
		\frac{p(1-q)}{q(1-p) + p(1-q)} \le \frac{1}{2} \iff Q \ge p
	.\] This implies that $\{X_n\}$ has probability 1 to eventually hit 0 iff $q \ge p$.

\end{proof}

\subsection{Null Recurrence and Positive Recurrence}

Let $S$ be a countable state space. 

\begin{defn}[Stationary distribution]
	We say that $\pi: S \to [0,1]$ is a stationary distribution for $\{X_n\}$ if $\sum\limits_{x \in S}^{} \pi_x = 1$ and for every $y \in S$, $\pi_y = \sum\limits_{x \in S}^{} p(x,y)$.
\end{defn}

The definition is the same as when $S$ was finite, in which case there was a unique stationary distribution if the Markov chain was irreducible and aperiodic. 

\begin{prop}
	Irreducible transient Markov chains do not have stationary distributions.
\end{prop}
\begin{proof}
	Assume that $\{X_n\}$ is transient with stationary distribution $\pi$. Then for any $x,y \in S$, \begin{equation*}\P\left[\{X_n\} \text{ visits $y$ finitely many times } \mid X_0 = x\right]  = 1.\end{equation*}\marginnote{This fact follows from the following: if $\{X_n\}$ visited $y$ infinitely many times, then $y$ would be recurrent, and since $\{X_n\}$ is irreducible, $\{X_n\}$ would be recurrent, a contradiction.} Hence, $p^n(x,y) = 0$ as $n \to \infty$. If $\pi$ is a stationary distribution, then for every $n \in \N$, \[
	\pi_y = \sum\limits_{x \in S}^{} \pi_x p^n(x,y) \to 0
\] as $n \to \infty$. However, since $y$ was arbitrary, this is a contradiction --- so transient Markov chains cannot have stationary distributions.
\end{proof}

\begin{defn}[Null recurrent, positive recurrent]
	Assume $\{X_n\}$ is irreducible and recurrent. We say that $\{X_n\}$ is null recurrent if $\lim\limits_{n \to \infty} p^n(x,y) = 0$ for every $x,y \in S$. $\{X_n\}$ is positive recurrent otherwise.

	Null recurrent Markov chains, like transient Markov chains, cannot have stationary distributions.
\end{defn}

\begin{prop}
	Assume that $\{X_n\}$ is irreducible. The following are equivalent:
	\begin{enumerate}
		\item $\{X_n\}$ is positive recurrent,
		\item $\{X_n\}$ has a stationary distribution,
		\item $\limsup\limits_{n \to \infty} p^n(x,y) > 0$ for every $x,y \in S$,
		\item If $T_x = \min \{n \ge 1:X_n = x\}$, then $\E\left[T_x \mid X_0 = x\right] < \infty$ for every $x \in S$.
	\end{enumerate}
	Furthermore, if $\{X_n\}$ is aperiodic and one of the above four conditions holds, then the stationary distribution $\pi$ is unique and \[
		\pi_y = \lim\limits_{n \to \infty} p^n(x,y) = \frac{1}{\E\left[T_y \mid X_0= y\right]}
	\] for any $x,y \in S$.
\end{prop}
\begin{proof}
	The proof is the same as in the finite state space case.
\end{proof}

If $\{X_n\}$ is irreducible, then for any $x \in S$, we have the following:
\begin{itemize}
	\item transient $\iff \sum\limits_{n=0}^{\infty} p^n(x,x) < \infty$,
	\item null recurrent $\iff \sum\limits_{n=0}^{\infty} p^n(x,x) = \infty$ and $\lim\limits_{n \to \infty} p^n(x,x) = 0$, and
	\item positive recurrent $\iff \limsup\limits_{n \to \infty} p^n(x,x) > 0$.
\end{itemize}

\texthl{It is likely for a question on the exam to be on categorizing a Markov chain as transient, null recurrent, or positive recurrent.} 

\begin{ex}
	Let $\{X_n\}$ be a biased random walk on $\N$ with partial reflected boundary, so
	\begin{align*}
		p(x,x-1) &= 1-p, \quad x\ge 1\\
		p(x,x+1) &= p, \quad x \ge 1 \\
		p(0,0) &= 1-p \\
		p(0,1) &= p
	.\end{align*}
	This Markov chain is irreducible on a countable state space. Since $p(0,0) > 0$, it is aperiodic. 

	(Case 1) When $p >\frac{1}{2}$, the Markov chain is transient, because there is a positive chance to never hit $0$ if $X_0 \ge 1$ (and hence it is not guaranteed to return to $0$ infinitely many times).

	(Case 2) Now suppose $p \le \frac{1}{2}$. We know that it is recurrent. To determine if it is null or positive recurrent, we try to find a stationary distribution and note \[
		\pi_y = \sum\limits_{x \in S}^{} \pi_x p(x,y)
	.\] For $y \ge 1$,
	\begin{align*}
		\pi_y &= p \pi_{y-1} + (1-p)\pi_{y+1} \\
		\pi_0 &= (1-p)\pi_0 + (1-p)\pi_1
	.\end{align*}

	(Case 2.1) We saw before that for $p < \frac{1}{2}$, the general solution is \[
		\pi_y = c_1 + c_2 \left( \frac{p}{1-p} \right)^{y}
	.\] Since we have the normalization constraint, we know $c_1 = 0$ and $1 = \sum\limits_{y=0}^{\infty} c_2 \left( \frac{p}{1-p} \right)^{y} =  \frac{c_2}{1-\frac{p}{1-p}}$. We take $c_2 = 1 - \frac{p}{1-p}$, so \[
		\pi_y = \left( 1- \frac{p}{1-p} \right)\left( \frac{p}{1-p} \right)^y
	\] characterizes a stationary distribution.

	(Case 2.2) When $p = \frac{1}{2}$, our solution is $\pi_y = c_1 + c_2 y$. Together with the normalization constraint, this implies that $c_1=c_2=0$, so $\pi_y = 0$ and there is no stationary distribution. Thus, the Markov chain is null recurrent.
\end{ex}

\begin{thm}
	Let $\{X_n\}$ be a random walk on $\Z^{d}$. View $\Z^d$ as a graph where $x,y \in \Z^d$ are joined by an edge iff $\left| x-y \right| = 1$. This random walk is null recurrent if $d \in \{1,2\}$ and transient if $d \ge 3$.
\end{thm}

\begin{lemma}
	(Stirling's Formula) Asymptotically, $n! \sim \sqrt{2\pi} n^{n+ \frac{1}{2}} e^{-n}$. (In the limit as $n\to \infty$, the ratio approaches $1$.
\end{lemma}

\begin{proof}
	($d=1$) Assume $X_0 = 0$. Since this is an irreducible Markov chain with countable state space, we just need to investigate recurrence/transience at 0. 

	If $n$ is odd, then $X_n = 0$, since $\Z$ is bipartite. We thus only need to consider even times. $X_{2n} = 0$ if and only if there are $n$ positive steps and $n$ negative steps among the first $2n$ steps of the Markov chain. The number of positive steps follows the binomial distribution, so

	\begin{align*}
		\P\left[X_{2n} = 0\right] &= \binom{2n}{n} \cdot \left( \frac{1}{2} \right)^{n} \left( \frac{1}{2} \right)^{n} \\
								  &= \frac{\left( 2n \right)!}{\left( n! \right)^2} \cdot \left( \frac{1}{2} \right)^{2n} 
	.\end{align*}
	Determining convergence from Stirling's formula,
	\begin{align*}
		\P\left[X_{2n} = 0\right] &\sim \frac{\sqrt{2\pi} \left( 2n \right)^{2n + \frac{1}{2}}e^{-2n}}{\left( \sqrt{2\pi} n^{n+\frac{1}{2}}e^{-n} \right)^2} \\
		&= \frac{1}{\sqrt{2\pi} } 2^{2n + \frac{1}{2}}n^{-\frac{1}{2}}2^{-2n} \\
		&=\frac{1}{\sqrt{\pi}n } \\
		\implies p^{2n}(0,0) &\to 0, \quad \sum\limits_{n=0}^{\infty} \P\left[X_{2n} = 0\right] = \infty
	.\end{align*}
	Hence, in dimension $1$, the random walk is null recurrent.

	($d\ge 2$, sketch) The random walk has $d$ components. The steps are i.i.d., so by the Law of Large Numbers, in $2n$ steps, about $\frac{2n}{d}$ steps should be in each direction. For $k=1,\ldots,d$, each step in the $k$th component is positive with probability $\frac{1}{2}$ and negative with probability $\frac{1}{2}$. By the $d=1$ case, \[
		\P\left[k\text{th component of $X_{2n}$ is } 0\right] = \frac{1}{\sqrt{\frac{\pi n}{d}} }
	.\] So,
	\begin{align*}
		\P\left[X_{2n} = 0\right] &= \left( \frac{1}{\sqrt{\frac{\pi n}{d}} } \right)^{d} \\
								  &= C \cdot n^{-\frac{d}{2}}\quad \text{for some $C \in \R$}
	.\end{align*}
	This goes to $0$ in the limit and is summable iff $d\ge 3$.
\end{proof}

\subsection{Branching Processes (Galton-Watson Processes)}

Consider an offspring distribution $\{p_k\}$ such that $p_k \ge 0$ represents the probability of having $k$ offspring and $\sum_{k=0}^{\infty} p_k =1$. Let $X_n$ denote the number of individuals in generation $n$. Each individual independently produces a number of offspring according to the distribution $\{p_k\}$, then dies.

Let $\xi_j$ be the number of offspring from individual $j$.
Note that $X_{n+1} = \sum\limits_{j=1}^{X_n} \xi_j$ where $\xi_j$ are conditionally independent given $X_n$ and for every $x \ge 0$, \[
	\P\left[\xi_j = k \mid X_n = x\right] = p_k
.\] Thus $\{X_n\}$ is a Markov chain.

What is the probability that our population goes extinct, i.e. $X_n = 0$ for some $n$? Define $a = \P\left[\exists n \ge 1 : X_n = 0\right]$. Put $\mu = \sum\limits_{k=0}^{\infty} k p_k$, which is the average of the offspring distributions. Intuitively, we are interested in if $\mu \ge 1$ or not.

Note that
\begin{align*}
	\E\left[X_{n+1} \mid X_n = m\right]  &= \E\left[\sum\limits_{j=1}^{m} \xi_j \mid X_n = m\right]  \\
	&= \sum\limits_{j=1}^{m} \E\left[\xi_j \mid X_n = m\right]  \\
	&= \mu m
.\end{align*}
Summing over all $m$, \[
	\E\left[X_{n+1}\right] = \mu \E\left[X_n\right] 
,\] so \[
	\E\left[X_n\right] = \mu^{n} \E\left[X_0\right] = \mu^{n}
\] if $X_0 = 1$. 

\begin{prop}
	If $\mu < 1$, then $a=1$. 
\end{prop}
\begin{proof}
	Assume $X_0 = 1$. Then $\P\left[X_n \ge 1\right] \le \E\left[X_n\right] = \mu^{n}\to 0$.

	$\P\left[X_n = 0\right] \to 1$ as $n\to \infty$, so $a = 1$.
\end{proof}

What if $\mu \ge 1$?
If $X_1 = k$, then descendants of the $k$ individuals are $k$ independent copies of the branching process. Then 
\begin{align*}
	\P\left[\text{extinction} \mid X_1 = k\right] &= a^{k}
.\end{align*}
Assuming that $X_0 = 1$, we have the following recursive equation:
\begin{align*}
	a= \P\left[\text{extinction}\right] &= \sum\limits_{k=0}^{\infty} \P\left[X_1 = k\right] \P\left[\text{extinction} \mid X_1 = k\right]  \\
										&= \sum\limits_{k=0}^{\infty} p_k a^k
.\end{align*}

\begin{defn}[Generating function]
	Let $Y$ be a random variable in $\N_0$. The generating function of $Y$ is $\phi = \phi_Y: [0,\infty] \to [0,\infty]$ defined by \[
		\phi(s) = \E\left[s^Y\right] = \sum\limits_{k=0}^{\infty} \P\left[Y=k\right] s^{k}
	.\] 
\end{defn}

Immediately, we see that $a = \phi(a)$ in the previous computation. We have the following basic properties of the generating function:
\begin{itemize}
	\item We allow $\phi(s) = \infty$ but $\phi(s) < \infty$ for all $s \in [0,1]$. 
	\item $\phi(1) = \sum\limits_{k=0}^{\infty} \P\left[Y= k \right] = 1$, \\ $\phi(0) = \P\left[Y=0\right] $
	\item $\phi'(s) = \sum\limits_{k=0}^{\infty} k \P\left[Y = k\right] s^{k-1}$, so $\phi'(1) = \E\left[Y\right] $
	\item If $Y_1,\ldots,Y_m$ are independent, then \begin{align*}
			\phi_{Y_1 + \ldots+Y_m}(s) &= \E\left[s^{Y_1 + \ldots + Y_n}\right] \\
			&= \prod_{j=1}^{m} \E\left[s^{Y_j}\right]   \\
			&= \prod_{j=1}^{m} \phi_{Y_j}(s) 
	.\end{align*}
\end{itemize}
\begin{prop}
Let $\phi$ be the generating function for $\{p_k\}$. Define $\phi^{(n)}$ to be $\underbrace{\phi \circ \cdots \circ \phi}_{n \text{ times}}$. Then $\phi_{X_n}(s) = \phi^{(n)}(s)$.
\end{prop}
\begin{proof}
	It suffices to show that $\phi_{X_{n+1}}(s) = \phi_{X_n} \left( \phi(s) \right)$.

	\begin{align*}
		\phi_{X_{n+1}}(s) &= \sum\limits_{k=1}^{\infty} \P\left[X_{n+1} = k\right] s^{k} \\
		&= \sum\limits_{k=0}^{\infty} \sum\limits_{j=0}^{\infty} \P\left[X_{n+1} = k \mid X_n = j\right] \P\left[X_n = j\right] s^{k} \\
		&= \sum\limits_{j=0}^{\infty} \left[\P\left[X_n = j\right] \sum\limits_{k=0}^{\infty} \P\left[X_{n+1} = k \mid X_n = j\right] s^{k}\right]
	.\end{align*}
	If $X_n = j$, then $X_{n+1} = \sum\limits_{i=1}^{j} \xi_j$ where $\{x_i\}$ are i.i.d according to $\{p_k\}$. So then $\sum\limits_{k=0}^{\infty} \P\left[X_{n+1} = k \mid X_n = j\right]s^{k}$ is the generating function of $\sum\limits_{i=1}^{j} \xi_i$, or $\left[ \phi(s) \right]^{j}$. Hence we can rewrite
	\begin{align*}
		\phi_{X_{n+1}}(s) &= \sum\limits_{j=0}^{\infty} \P\left[X_n = j \right] \left[ \phi(s) \right]^{j} \\
						  &= \phi_{X_n}(\phi(s))
	.\end{align*}
\end{proof}

\begin{prop}
	Assume that $0 < p_0 < 1$. Then $a$ is the smallest positive solution of $\phi(s) = s$.
\end{prop}
\begin{proof}
	We already know that $\phi(a) = a$. First we need to check that there exists a smallest positive solution. Note that $\phi(s) - s$ is continuous on $[0,1]$. Hence $\{s \in [0,1]: \phi(s) = s\}$ is closed.

	Since $\phi(0) = p_0 > 0$ and $\phi(1) = 1$, $0$ is not in the set and $1$ is in the set. This set has an element greater than $0$, is closed, and does not contain $0$, so there exists a smallest positive $s_0 > 0$ in $\{s in [0,1] : \phi(s) = s\}$. 

	We claim that if $s_0$ is the smallest positive solution, then \[
		\phi_{X_n}(0) < s_0
	.\] Considering the base case, $\phi_{X_0}(s) = s$. In particular, $\phi_{X_0}(0) = 0 < s_0$. Now assume $n \ge 0$ and $\phi_{X_n}(0) < s_0$. Then 
	\begin{align*}
		\phi_{X_{n+1}}(0) &= \phi\left( \underbrace{\phi_{X_n}(0)}_{< s_0} \right) \\
						  &< \phi(s_0) = s_0
	,\end{align*}
	where the inequality follows from $\phi$ being increasing.

	Since $\phi_{X_n}(0) = \P\left[X_n = 0\right] < s_0$ and $\phi(a) = a = \lim\limits_{n \to \infty} \P\left[X_n = 0\right]  \le s_0$, $a$ must be the smallest positive solution, so $a = s_0$.
\end{proof}
\begin{prop}
	Assume $0 < p_0 < 1$. If $\mu > 1$, then $a < 1$. If $\mu \le 1$, then $a = 1$.
\end{prop}
\begin{proof}
	We saw earlier that $a=1$ if $\mu < 1$. So we only need to consider when $\mu \ge 1$.

	Begin by looking at the second derivative of $\phi$:
	\begin{align*}
		\phi''(s) &= \sum\limits_{k=2}^{\infty} k(k-1)p_k s^{k-2}
	.\end{align*}
	If $p_0 > 0$ and $\mu \ge 1$, then $p_k > 0$ for some $k \ge 2$. Thus $\phi''(s) > 0$ for all $s \in (0,1)$, so $\phi'$ is increasing on the interval. Furthermore, recall that for any $s \in (0,1)$, $\phi(1) = 1$ and $\phi'(1) = \mu$. 

	(Case 1: $\mu =1$) By the FTC, 
	\begin{align*}
		1- \phi(s) &= \int_{s}^{1} \phi'(t)dt \\
		&< 1-s \quad \text{ since } \phi'(1) = 1, \phi' \text{ increasing} \\
		\implies \phi(s) > s \quad \forall  s \in [0,1)
	.\end{align*}
	No $s$ satisfies this, so $a=1$.

	(Case 2: $\mu > 1$) Now $\phi'(1) > 1$ and $\phi(1) = 1$. Doing a Taylor expansion at $s = 1$, 
	\begin{align*}
		\phi(1-\epsilon) &= 1 -\phi'(1)\epsilon + o(\epsilon) \\
		\implies \exists \epsilon \in (0,1) \text{s.t. } \phi(1-\epsilon) &< 1- \epsilon
	.\end{align*}
	$\phi(0) = p_0 > 0$. Since $\phi(s)$ is continuous, there exists $t \in (0,1-\epsilon)$ such that $\phi(t) = t$. Hence $a \le 1- \epsilon$.
\end{proof}

\begin{defn}[Critical branching process]
	$\{X_n\}$ is
	\begin{itemize}
		\item supercritical if $\mu > 1$ (and hence $a < 1$)
		\item critical if $\mu = 1$ (and hence $a=1$)
		\item subcritical if $\mu < 1$ (and hence $a=1$)
	\end{itemize}
\end{defn}

\begin{ex}
	Let $p_0 = \frac{1}{10}, p_1 = \frac{3}{5}, p_2 = \frac{3}{10}$. Then \[
		\mu = \frac{3}{5} + \frac{3}{10} * 2 = \frac{6}{5} > 1
	,\] so this is a supercritical branching process.
	To compute the extinction probability, we solve
	\begin{align*}
		\phi(s) &= s \\
		\frac{1}{10}s^{0} + \frac{3}{5} s^{1} + \frac{3}{10}s^{2} &=  s \\
		\frac{1}{10} + \frac{3}{5}s + \frac{3}{10}s^2 &= 0 \\
		\implies s = 1, \frac{1}{3}
	.\end{align*}
	The extinction probability is $\frac{1}{3}$.
\end{ex}

\section{Continuous-time Markov Processes}

\subsection{Poisson Processes}

Consider a stochastic process with continuous time and discrete space.

\marginnote{We might want to model the number of phone calls $X_t$ received at or before time $t \ge 0$,
where there are an average of $\lambda > 0$ calls per hour and the number of calls during disjoin time intervals are independent. In this case, $Y$ models the numbers of calls in one hour.}
\begin{defn}[Poisson distribution]
	A random variable $Y$ has the Poisson distribution with mean $\lambda$ if \[
		\P\left[Y = k\right] = \frac{\lambda^{k}e^{-\lambda}}{k!}
	\] for every $k=0,1,2,\ldots$.
\end{defn}

\begin{defn}[Poisson process]
	The Poisson process with rate $\lambda$ is the continuous-time stochastic process $\{X_t\}_{t\ge 0}$ such that $X_0 =0$ and for any times \[
		 0 \le s_1 \le t_1\le s_2 \le t_2\le \ldots\le s_k \le t_k
	,\] the random variables \[
		X_{t_j} - X_{s_j}, \quad j = 1,\ldots,k
	\] are independent, and each has a Poisson distribution with mean $\lambda(t_j-s_j)$.
\end{defn}
\marginnote{We can describe a Poisson process in terms of the arrival times: Let $\tau_0 = 0$ and let \[
	\tau_j = \inf \{t \ge 0 : X_t = j\}, \quad \forall  j=1,2,\ldots
\] be the time that the $j$th call comes. We have that \[
	X_t = \max \{j: \tau_j \le t\}
.\] }


\begin{defn}[Exponential distribution]
	A random variable $T \in [0,\infty)$ has the exponential distribution with parameter $\lambda$ if for each $t\ge 0$, \[
		\P\left[T \ge t\right]  = e^{-\lambda t}
	.\] Integrating, we find that the mean of the exponential distribution is $\frac{1}{\lambda}.$
\end{defn}

\begin{prop}
	The arrival times $\tau_j - \tau_{j-1}$ for $j=1,2,\ldots$ for a Poisson process with rate $\lambda$ are i.i.d. according to the exponential distribution with parameter $\lambda$. 
\end{prop}
\begin{proof}
	Note that \[
		\P\left[\tau_1 > t\right] = \P\left[X_t = 0\right] = e^{-\lambda t}
	,\] so $\tau_1 \sim Exp(\lambda)$. By the independent increments property, for any $t \ge 0$, $\{X_{s+t} - X_t\}$ is a Poisson process independent from $\{X_s\}_{s\le t}$. The time $\tau_1$ is a stopping time for $\{X_s\}_{s\le t}$; i.e., $\{\tau_1 \le t\}$ is determined by $\{X_s\}_{s\le t}$.

	Using a version of the strong Markov property for the Poisson process, we get that $\{X_{s+\tau_1}-X_{\tau_1}\}_{s\ge 0}$ is a Poisson process independent from $\{X_s\}_{s \le \tau_1}$. In particular, $\tau_2-\tau_1$ is independent from $\tau_1$ and has the same distribution as $\tau_1$. We can iterate on this process to get that all increments $\tau_j - \tau_{j-1}$ are i.i.d.
\end{proof}

\begin{rmk}
	We can construct a Poisson process as follows. Let $\{T_j\}_{j\ge 1}$ be i.i.d. exponential random variables with parameter $\lambda$ and let \[
		X_t = \max \{k \ge 0 : \sum\limits_{j=1}^{k} T_j \le t\}
	.\] Then $\{X_t\}$ is a Poisson process.
\end{rmk}

\begin{ex}
	Suppose at a certain bus station, buses arrive according to a Poisson process with a rate of 2 buses per hour.
	If we wait for two hours, what is the probability that we have not seen a bus yet?
	\newthought{Solution.} The number in the first two hours is found from the exponential distribution with $\P\left[T \ge 2\right]$ at $\lambda=2$. The probability that there are zero buses is $e^{-(2)(2)} \approx 0.0183$.
\end{ex}

\begin{ex}
	Given that no bus arrived in the first two hours, what is the conditional probability that a bus will arrive in the next hour?

	\newthought{Solution.} The number of buses between times 2 and 3 is independent from the number of buses before time 2, so the number of buses between times 2 and 3 follows $Poi(2)$; the probability that this number is nonzero is $1- \frac{2^0 e^{-2}}{0!} \approx 0.864$.
\end{ex}

\begin{ex}
	What is the expected time that the second bus arrives?

	\newthought{Solution.} The second arrival time is the sum of two exponential random variables with parameter 2 (mean $\frac{1}{2}$), so its mean is 1.
\end{ex}

\subsection{More on the Exponential Distribution}

We have the following properties of the exponential distribution:
\begin{itemize}
	\item Memoryless property: For any $t\ge 0$, the conditional distribution of $T-t$ given $\{T \ge t\}$ is also $Exp(\lambda)$: \[
		\P\left[T \ge s+ t \mid T \ge t\right]  = \frac{e^{-\lambda(s+t)}}{e^{-\lambda t}} = e^{-\lambda s}
	.\] 
	\item Minimum property: Let $T_1,\ldots,T_n$ be independent exponential random variables with parameters $\lambda_1,\ldots,\lambda_n$. Then \[
		\min \{T_1,\ldots,T_n\} \sim Exp(\lambda_1+\ldots+\lambda_n)
	\] and \[
		\P\left[T_j = \min \{T_1,\ldots,T_n\}\right] = \frac{\lambda_j}{\lambda_1+\ldots+\lambda_n}
	.\] 
	\begin{itemize}
		\item To get the distribution of $\min \{T_1,\ldots,T_n\}$, we can compute \[
		\P\left[\min \{T_1,\ldots,T_n\} \ge t\right] = \prod_{k=1}^{n} \P\left[T_k \ge t\right] = \prod_{k=1}^{n} e^{-\lambda_k t} = e^{-(\lambda_1+\ldots+\lambda_n)t}
	.\] 
		\item To get the other claim, we assume that $j=1$ by relabeling, then compute
			\begin{equation}
				\begin{aligned}
					\P\left[T_1 = \min \{T_1,\ldots,T_n\}\right] &= \P\left[T_k > T_1, \forall k=2,\ldots,n\right]  \\
					&= \E\left[\P\left[T_k > T_1, \forall k=2,\ldots,n \mid T_1\right] \right]  \\
					&= \E\left[e^{-\left( \lambda_2+\ldots+\lambda_n \right)T_1}\right] \\
					&= \int_{0}^{\infty} e^{-(\lambda_1+\ldots+\lambda_n)t}\lambda_1 dt  \\
					&= \frac{\lambda_1}{\lambda_1+\ldots+\lambda_n}
				\end{aligned}
			\end{equation}
	\end{itemize}
\end{itemize}

\subsection{Continuous Time Markov Chains}

Let $S$ be a finite state space. We will define a stochastic process $\{X_t\}_{t\ge 0}$ taking values in $S$.\marginnote{The Markov chain jumps from one state to another, similarly to how the Poisson process jumps from one integer to the next.} Define rates $\alpha(x,y) \ge 0$ for distinct $x,y \in S$ that specify how frequently we jump from $x$ to $y$. 

Assume that we start in state $x$. For $y \in S$ with $\alpha(x,y) \ne 0$, let $T_y$ be an exponential rnadom variable with parameter $\alpha(x,y)$. We consider the $T_y$s for different $y$ to be independent.\marginnote{$T_y$ can be thought of as an alarm clock that rings at time $T_y$.} Put $T = \min_y t_y$ so that at time $T$ we move from state $x$ to state $y$ (and $y$ was chosen such that $T_y$ is minimal). 

By the minimum property of the exponential distribution, \[
	T \sim Exp(\alpha(x)), \quad \alpha(x) = \sum\limits_{\substack{y \in S \\ y \ne x}}^{} \alpha(x,y)
.\] Also, \[
	\P\left[X_T = y\right] = \frac{\alpha(x,y)}{\alpha(x)}
.\]\marginnote[-3\baselineskip]{These properties follow from the minimum property in the previous section on the exponential distribution.}

After jumping from $x$ to $X_T$, choose the next jump redefining the $T_y$s as new exponential random variables with parameters $\alpha(X_T, z)$ for each $z \in S$. If for some $x \in S$, we have that $\alpha(x,y) = 0$ for every $y \in S$, then the Markov chain stays at $x$ forever if it reaches $x$.

\subsection{Relating Continuous and Discrete Time Markov Chains}

The midterm will cover material through the end of this section.

Consider a finite state space $S$ and a continuous time Markov chain $\{X_t\}_{t\ge 0}$ with rates $\alpha(x,y)$ defined for $x\ne y$ and value $\alpha(x) = \sum\limits_{\substack{y \in S \\ y \ne x}}^{} \alpha(x,y)$.

If $X_0 = x \in S$, let $T \sim Exp(\alpha(x))$ be the time at which the Markov chain jumps to $y \ne x$. We have that $\P\left[X_T = y\right] = \frac{\alpha(x,y)}{\alpha(x)}$. The Markov chain continues iteratively, redefining $T$. 

\begin{prop}
	\marginnote{Note that $t$ is likely not a jump time; this version of the discrete Markov chain property holds for all $t$.}Let $t \ge 0$ and $x \in S$. Then the conditional distribution of $\{X_{t+s}\}_{s\ge 0}$ given $\{X_t = x\}$ and everything before time $t$ is the same as the original Markov chain $\{X_s\}$ started at $x$.\marginnote{The proof relies on the memoryless property of the exponential distribution, but is too advanced for this course.}
\end{prop}

\newthought{How are continuous} Markov chains related to discrete Markov chains? Let $\tau_0 = 0$ and $\tau_1,\tau_2,\ldots$ be the following consecutive jump times. If we condition on $\{X_{\tau_0},\ldots,X_{\tau_{n-1}} = x\}$, then conditional probability that $X_{\tau_n} = y$ is given by \[
	\frac{\alpha(x,y)}{\alpha(x)}
.\] 

This implies that we can turn continuous Markov chains into discrete ones by just examining the values at jump times. \marginnote{Note that this process implies that $p(x,x) = 0$, as the continuous Markov chain cannot jump to the same state that it is in.}Putting $\tilde{X}_n = X_{\tau_n}$ for $n \in \N$, then $\{\tilde{X}_n\}$ is a discrete time Markov chain with transition probabilities \[
	p(x,y) = \frac{\alpha(x,y)}{\alpha(x)}, \quad y \ne x
.\] 

This is called the \emph{associated discrete time Markov chain}. You cannot recover the continuous Markov chain from the associate discrete chain because the discrete version does not encode the waiting times $\tau_n - \tau_{n-1}$.

\begin{defn}[Time-$t$ transition probability]
	Define the time-$t$ transition probability as \[
		p_t(x,y) = \P\left[X_t = y \mid X_0= x\right] , \quad t \ge 0, x,y \in S
	.\] This is analogous to the $n$-step transition probabilities. 
\end{defn}
WLOG, we can assume that $S = \{1,\ldots,N\}$.
This allows us to define the following.

\begin{defn}[Time-$t$ transition matrix]
	Define the $N \times N$ matrix \[
		P_t = \left( p_t(x,y) \right)_{x,y = 1,\ldots,N}
	.\] 
\end{defn}

\newthought{Placeholder.} For $x \ne y$, we can see that\marginnote[\baselineskip]{Put $T_y \sim Exp(\alpha(x,y))$.}\marginnote[2\baselineskip]{The $o(\epsilon)$ term comes from the fact that the second minimal $T_z$ and so on can also be less than $\epsilon$.}
\begin{align*}
	\P\left[X_{t+\epsilon} = y \mid X_t = x\right] &= \P\left[X_{\epsilon} = y \mid X_0 = x\right] \\
	&= \P\left[T_y < \epsilon \mid X_0 = x\right] + o(\epsilon)  \\
	&= 1- e^{-\alpha(x,y) \epsilon} + o(\epsilon) \\
	&= \alpha(x,y)\epsilon + o(\epsilon) \quad \text{by Taylor expansion}\\
	\text{Hence,} \quad \P\left[X_{t+\epsilon} \ne x \mid X_t = x\right] &= \sum\limits_{y \ne x}^{} \P\left[X_{t+\epsilon} = y \mid X_t = x\right] \\
																		 &= \sum\limits_{y\ne x}^{} \alpha(x,y)\epsilon + o(\epsilon) \\
																		 &= \alpha(x)\epsilon + o(\epsilon)
																		 \implies \frac{1}{\epsilon}\P\left[X_{t+\epsilon} \ne x \mid X_t= x\right] = \alpha(x) + o(\epsilon)
.\end{align*}
The last expression can be plugged in to the last line of the following expansion to find

\begin{align*}
	\frac{d}{dt} p_t(x,y) &= \lim\limits_{e \to 0} \frac{1}{\epsilon} \left( \P\left[X_{t+\epsilon} = y \mid X_0 = x\right] - \P\left[X_t = y \mid X_0 = x\right]  \right) \\
						  &= \lim\limits_{\epsilon \to 0} \frac{1}{\epsilon}\left( \P\left[X_{t+\epsilon} = y, X_t = y \mid X_0 = x\right] + \P\left[X_{t + \epsilon} = y, X_t \ne y \mid X_0=x\right] - \right. \\
						  & \qquad\left. \P\left[X_{t + \epsilon} = y, X_t = y \mid X_0 = x\right] + \P\left[X_{t+\epsilon} \ne y, X_t = y \mid X_0 = x\right]  \right) \\
	&= \lim\limits_{\epsilon \to 0} \frac{1}{\epsilon}\left( \sum\limits_{z \ne y}^{} \P\left[X_{t + \epsilon} = y \mid X_t = z\right] p_t(x,z) - \P\left[X_{t+\epsilon} \ne y, X_t = y \mid X_0 = x\right] \right) \\
	&= \sum\limits_{z \ne y}^{\infty}  \alpha(z,y)p_t(x,z) - \alpha(y)p_t(x,y)
.\end{align*}

\begin{defn}[Infinitesimal generator]
	The infinitesimal generator is the $N\times N$ matrix $A$ where \[
		A_{x,y} = \begin{cases}
			\alpha(x,y), \quad y \ne x \\
			-\alpha(x), \quad y = x
		\end{cases}
	.\] It looks like \[
		A = \begin{pmatrix}
		    -\alpha(1) & -\alpha(1,2) & \cdots & \alpha(1,N) \\
			\alpha(2,1) & \ddots & & \vdots \\
			\alpha(N,1) & \cdots & & -\alpha(N)
		\end{pmatrix} 
	.\] 
\end{defn}

With the initial condition matrix given by the entries $p_0(x,x) = 1$ and $p_0(x,y) = 0$ for every $y \ne x$ (since the Markov chain cannot jump right where it starts), we get an ODE for the matrix-valued function $p_t$, where 
\begin{align*}
	\frac{d}{dt} P_t &= P_t A \\
	P_0 = I
.\end{align*}
This implies that 
\begin{align*}
	P_t &= e^{t A} = \sum\limits_{n=0}^{\infty} \frac{t^n A^n}{n!}, \\
	\text{which we verify with } \frac{d}{dt} e^{tA} &= e^{tA} A = A e^{tA} \\
	e^{0 A} &= I
.\end{align*}

If $A$ is diagonalizable such that $A = QDQ^{-1}$, then $e^{tA} = Q e^{t D} Q^{-1}$. Hence, if the matrix $A$ as defined above is diagonalizable, we can easily compute the time-$t$ transition matrix.

\begin{ex}
	Suppose $S = \{0,1\}$, $\alpha(0,1) = 2$, and $\alpha(1,0)$ = 3. Then $\alpha(0) = 2, \alpha(1) = 3$, so \[
		A = \begin{pmatrix}
		-2 & 2 \\
		3 & -3
		\end{pmatrix} 
	.\] We know that the transition matrix takes on the form $P_t = e^{tA}$. 
	
	We find the eigenvalues by solving \[
		0 = det(A - \lambda I) = (-2-\lambda)(-3-\lambda) - 6
	,\] which implies that $\lambda_1 = -5$, $\lambda_2 = 0$, so by putting $Av = \lambda v$, we find $v_{-5} = (-2,3)$ and $v_0 = (1,1)$.

	Writing $A = QDQ^{-1}$, where \[
		Q = \begin{pmatrix}
		-2 & 1 \\
		3 & 1
		\end{pmatrix} \quad D = \begin{pmatrix}
		5 & 0 \\
		0 & 0
		\end{pmatrix} 
	.\] 
	Hence \[
		P_t = Qd^{t D} Q^{-1} = \frac{1}{5} \begin{pmatrix}
		    3+2e^{-5 t} & 2-2e^{-5 t} \\
			3-3e^{-5t} & 2+3e^{-5t}
		\end{pmatrix} 
	.\] 
	We can read off individual transition probabilities: \[
		p_t(0,1) = \frac{1}{5} \left( 2-2e^{-5t} \right)
	.\] 
\end{ex}

\subsection{Stationary Disributions}

\begin{defn}[Irreducible]
	We say that the continuous-time Markov chain $\{X_t\}$ is irreducible if we can get from any state to any other state.
\end{defn}
As opposed to the discrete case, irreducibility implies that $p_t(x,y) > 0$ for all $t>0$ and $x,y \in S$. In the discrete case, it may not have been possible to go from $x$ to $y$ in 1 step, but in the continuous case, our jump times are exponentially distributed and can be arbitrarily small.

There is no notion of periodicity in continuous time Markov chains.

\begin{defn}[Stationary distribution]
	The function $\pi: S \to [0,1]$ with $\sum\limits_{x \in S}^{} \pi_x = 1$ is a stationary distribution for $\{X_t\}$ if whenever $X_0 \sim \pi$, then $X_t \sim \pi$ for all $t \ge 0$.

	Note that $\pi$ is stationary iff \[
		\pi P_t = \pi, \quad \forall t \ge 0
	,\] which is equivalent to
	\begin{align*}
		& \quad \frac{d}{dt} \pi P_{t} = 0 \\
		& \iff \pi \frac{d}{dt}e^{tA} = 0 \\
							   &\iff \pi A e^{tA} = 0 \\
							   &\iff \pi A = 0
	,\end{align*}
	since $A$ is invertible. Hence we get that $\pi$ is stationary iff $\pi A = 0$ and $\sum\limits_{x \in X}^{} \pi_x = 1$.
\end{defn}

\begin{prop}
	If $\{X_t\}$ is irreducible, then there exists a unique stationary distribution $\pi$ for $\{X_t\}$. Furthermore, \[
		\lim\limits_{t \to \infty} p_t(x,y) = \pi_y, \quad \forall x,y \in S
	.\] 
\end{prop}

\begin{ex}
	Given $A = \begin{pmatrix}
	-2 & 2 \\
	3 & -3
	\end{pmatrix}$, we can solve
	\begin{align*}
		(\pi_0 \quad \pi_1) \begin{pmatrix}
		-2 & 2 \\
		3 & -3
	\end{pmatrix} &= (0 \quad 0) \\
	\pi_0 + p_1 &=  1
	\end{align*}
	to find that $\pi_0 = \frac{3}{5}$, $\pi_1 = \frac{2}{5}.$
\end{ex}

\subsection{Summary: Continuous vs. Discrete}

In discrete time, we have
\begin{itemize}
	\item transition probabilities $p(x,y)$ encoded in the transition matrix $P$
	\item $n$-step transition probabilities $p^n(x,y) = \left( P^n \right)_{x,y}$
	\item the stationary distribution $\pi$ solves $\pi P = \pi$, $\sum\limits_{x \in S}^{} \pi_x =1$
.\end{itemize}
In continuous time, we have
\begin{itemize}
	\item transition rates $\alpha(x,y)$ that are encoded in the infinitesimal general matrix $A$
	\item time-$t$ transition probabilities $p_t(x,y) = \left( P_t \right)_{x,y}$, where $P_t = e^{t A}$
	\item the stationary distribution $\pi$ solves $\pi A = 0$, $\sum\limits_{x \in S}^{} \pi_x = 1$
.\end{itemize}

\section{Martingales}
\subsection{Conditional Expectation: Introduction}

Let $X$ be a random value in $\R$ and $Y$ be a random variable.

For any $y \in T \subseteq \R$, where $T$ and $S$ are countable, we have that
\begin{align*}
	\E\left[X \mid Y = y\right] &= \sum\limits_{s \in S}^{} s \P\left[X = s \mid Y = y\right]  \\
								&= \frac{\sum\limits_{s \in S}^{} s f(s,y)}{\sum\limits_{s \in S}^{} f(s,y)}
.\end{align*}

Hence, we define conditional expectation for random variables $X,Y$ with countable support, the conditional expectation of $X$ with respect to $Y$ is \[
	\E\left[X \mid Y\right] = \frac{\sum\limits_{s\in S}^{} s f(s,Y)}{\sum\limits_{s \in S}^{} f(s,Y)}
.\] \marginnote{Note that $\E\left[X \mid Y\right] $ is a random variable that is some function of $Y$.}

\begin{ex}
	Let $X,Y$ be random variables in $\R$. Assume there exists a joint density $f(x,y)$. Then \[
		\P\left[X \in A, Y \in B\right] = \int_{A}^{} \int_{B}^{} f(x,y) dxdy  
	.\] 
\end{ex}

If $X,Y$ are continuous random variables, then the conditional density of $X$ given $Y$ is \[
	f(x \mid y) = \frac{f(x,y)}{\int_{\R}^{} f(u,y)du }
.\] 

We can get the conditional expectation by integrating: \[
	\E\left[X \mid Y\right] = \frac{\int_{\R}^{} xf(x,Y)dx }{\int_{\R}^{} f(u,Y)du }
.\] 

We want a singular definition of conditional expectation to satisfy both the discrete and continuous cases.

\begin{defn}[Conditional expectation]
	Let $X,Y$ be random variables such that $X \in \R$ and $\E\left[\left| X \right|\right] < \infty$. The conditional expectation $\E\left[X \mid Y\right] $ is the unique random variable in $\R$ such that:
	\begin{enumerate}[(i)]
		\item $\E\left[X \mid Y\right] $ is a function of $Y$, and 
		\item If $F(Y)$ is a function of $Y$ that takes values in $\R$ with $\E\left[\left| F(X) \right|\right] < \infty$, then \[
			\E\left[X F(Y)\right] = \E\left[\E\left[X\mid Y\right] F(Y)\right] 
		.\] 
	\end{enumerate}
\end{defn}

\begin{thm}
	The conditional expectation exists and is unique.
\end{thm}
\begin{proof}
	Take Graduate Analysis III with Gwynne.
\end{proof}

\begin{defn}[Conditional probability]
	The conditional probability of an event $E$ given $Y$ is \[
		\P\left[E \mid Y\right] = \E\left[\mathds{1}_E \mid Y\right] 
	.\] 
\end{defn}

\begin{ex}
	Suppose $S,T$ are countable, $S \subset \R$, and $f(s,y) = \P\left[X = s, Y = y\right] $.
	We defined the conditional expectation \[
		\E\left[X \mid Y\right]  = \frac{\sum\limits_{s\in S}^{} sf(s,Y)}{\sum\limits_{s\in S}^{} f(s,Y)}
	.\] Does this satisfy the abstract definition?
	\begin{itemize}
		\item It is a function of $Y$.
		\item Let $f(Y)$ be a function of $Y$ taking values in $\R$ with $\E\left[\left| F(Y) \right|\right] < \infty$. Then
			\begin{align*}
				\E\left[\E\left[X\mid y\right] F(Y)\right]  &= \sum\limits_{u \in S}^{} \sum\limits_{v \in T} f(u,v) F(v) \frac{\sum\limits_{s \in S}^{} s f(s,v)}{\sum\limits_{s\in S}^{} f(s,v)}\\
				&= \sum\limits_{v \in T}^{}  \left( \sum\limits_{u\in S}^{} f(y,v) \right) F(v) \frac{\sum\limits_{s \in S}^{} s f(s,v)}{\sum\limits_{s \in S}^{} f(s,v)}\\
				&= \sum\limits_{v \in T}^{} \sum\limits_{s \in S}^{} F(v) s f(s,v) \\
				&= \E\left[X F(Y)\right] 
			.\end{align*}
	\end{itemize}
\end{ex}

\subsection{Properties of conditional expectation}

Suppose all random variables $X$ taking values in $\R$ satisfy $\E\left[\left| X \right|\right] < \infty$. Given a sequence of random variables $\{Y_n\}_{n \ge 0}$, we can define $\E\left[X \mid Y_0,\ldots,Y_n\right] $ by taking $Y = \left( Y_0,\ldots,Y_n \right)$.

\begin{defn}[$\mathscr{F}_n$-measurable]
	Let $\mathscr{F}_n$ be the information contained in $Y_0,\ldots,Y_n$. More precisely, let \[
		\E\left[X \mid \mathscr{F}_n\right] = \E\left[X \mid Y_0,\ldots,Y_n\right] 
	.\] We think of $n$ as time and $\mathscr{F}_n$ representing the information available to us at time $n$. A random variable $X$ is $\mathscr{F}_n$-measurable if $X$ is a function of $Y_0,\ldots,Y_n$.\marginnote{$\mathscr{F}_n$ is a $\sigma$-algebra generated by $Y_0,\ldots,Y_n$.}
\end{defn}

\begin{prop}
	$\E\left[X \mid \mathscr{F}_n\right] $ is the unique random variable taking values in $\R$ satisfying 
	\begin{enumerate}[(i)]
		\item $\E\left[X \mid \mathscr{F}_n\right] $ is $\mathscr{F}_n$-measurable
		\item If $Z$ is an $\mathscr{F}_n$-measurable random variable in $\R$, then \[
			\E\left[XZ\right] = \E\left[\E\left[X \mid \mathscr{F}_n\right] Z\right] 
		.\] 
	\end{enumerate}
\end{prop}

\begin{prop}
	(Linearity) If $X_1,X_2$ are random variables in $\R$ and $a,b \in \R$ are constants, then \[
		\E\left[aX_1 + bX_2 \mid \mathscr{F}_n\right] = a\E\left[X_1 \mid \mathscr{F}_n\right] + b \E\left[X_2 \mid \mathscr{F}_n\right] 
	.\] 
\end{prop}

\begin{proof}
	Let \[W = a \E\left[X_1 \mid \mathscr{F}_n\right] + b \E\left[X_2 \mid \mathscr{F}_n\right]. \] We want to check that $W$ satisfies the definition of of $\E\left[aX_1 + bX_2 \mid \mathscr{F}_n\right] $, since by uniqueness of conditional expectation, we will know they're equivalent.

	The two terms in $W$ are $\mathscr{F}_n$-measurable, as they are functions of $Y_0,\ldots,Y_n$ by definition. Hence $W$ is a function of $Y_0,\ldots,Y_n$ and thus $\mathscr{F}_n$-measurable.
	Furthermore, we can compute
	\begin{align*}
		\E\left[WZ\right]  &= a\E\left[\E\left[X_1 \mid \mathscr{F}_n\right] Z\right] + b\E\left[\E\left[X_1 \mid \mathscr{F}_n\right] Z\right] \\
		&= a \E\left[X_1Z\right] + b\E\left[X_2Z\right]  \\
		&= \E\left[\left( aX_1 + bX_2 \right)Z\right] 
	.\end{align*}
\end{proof}

\begin{prop}
	If $X$ is $\mathscr{F}_n$-measurable, then $\E\left[X \mid \mathscr{F}_n\right] = X$.
\end{prop}

\begin{proof}
	$X$ is $\mathscr{F}_n$ measurable by assumption. If $Z$ is $\mathscr{F}_n$ measurable, then
	\begin{align*}
		\E\left[XZ\right] = \E\left[XZ\right] ,
	\end{align*} so we are done.
\end{proof}

\begin{prop}
	If $X$ is independent from $\mathscr{F}_n$, then \[
		\E\left[X \mid \mathscr{F}_n\right]  = \E\left[X\right] 
	.\] 
\end{prop}
\begin{proof}
	$\E\left[X\right]$ is $\mathscr{F}_n$-measurable because it is a nonrandom constant.
	Now suppose $Z$ is $\mathscr{F}_n$-measurable. We see that
	\begin{align*}
		\E\left[XZ\right]  &= \E\left[X\right] \E\left[Z\right] \quad \text{by independence} \\
						   &= \E\left[\E\left[X \right] Z\right] \quad \text{\texthl{by linearity (?)}}
	.\end{align*}
\end{proof}

\begin{prop}
	If $Z$ is $\mathscr{F}_n$-measurable, then \[
		\E\left[ZX \mid \mathscr{F}_n\right] = Z \E\left[X \mid \mathscr{F}_n\right] 
	.\] 
\end{prop}

\begin{proof}
	Let $W = Z \E\left[X \mid \mathscr{F}_n\right] $.
	$W$ is the product of $Z$ and $\E\left[X \mid \mathscr{F}_n\right]$, which are both $\mathscr{F}_n$-measurable by definition.
	For the second property, let $Z'$ be another $\mathscr{F}_n$-measurable random variable. Then 
	\begin{align*}
		\E\left[W Z'\right] &= \E\left[\underbrace{Z' Z}_{\mathscr{F}_n\text{-measurable}} \E\left[X \mid \mathscr{F}_n\right] \right]  \\
		&= \E\left[Z' \cdot ZX\right]
	.\end{align*}
\end{proof}

\begin{prop}
	\[
		\E\left[\E\left[X \mid \mathscr{F}_n\right] \right] = \E\left[X\right] 
	.\] 		
\end{prop}
\begin{proof}
	Take $Z = 1$ in the definition of conditional expectation.
\end{proof}

\begin{prop}
	(Tower property) If $m \le n$, then \[
		\E\left[\E\left[X \mid \mathscr{F}_n\right] \mid \mathscr{F}_n\right] = \E\left[X \mid \mathscr{F}_m\right] 
	.\] 
\end{prop}
\begin{proof}
	We know by definition that $
		\E\left[\E\left[X \mid \mathscr{F}_n\right] \mid \mathscr{F}_m\right] 
	$ is $\mathscr{F}_m$-measurable.

	Furthermore, if $Z$ is $\mathscr{F}_m$-measurable, then
	\begin{align*}
	\E\left[Z \E\left[\E\left[X \mid \mathscr{F}_n\right] \mid \mathscr{F}_m\right] \right] &= \E\left[Z \E\left[X \mid \mathscr{F}_n\right] \right] \end{align*} by the defining property of $\E\left[X \mid \mathscr{F}_n\right] $, which leads us to equality with $\E\left[ZX\right] $.
\end{proof}

\begin{prop}
	\[
		\left| \E\left[X \mid \mathscr{F}_n\right]  \right| \le \E\left[\left| X \right|\mid \mathscr{F}_n\right] 
	.\] 	
\end{prop}

\begin{proof}
	Let $Z$ be a $\mathscr{F}_n$-measurable random variable in $[0,\infty)$. Then
	\begin{align*}
		\E\left[Z \E\left[X \mid \mathscr{F}_n\right] \right] &= \E\left[ZX\right]  \\
		&\le \E\left[Z \left| X \right|\right]  \\
		&= \E\left[Z \E\left[\left| x \right| \mid \mathscr{F}_n\right] \right]
	.\end{align*}
	This holds for any nonnegative random variable of $Z$, so \[
		\E\left[X \mid \mathscr{F}_n\right] \le \E\left[\left| X \right|\ \mid \mathscr{F}_n\right] 
	.\] By a similar argument, \[
		-\E\left[X \mid \mathscr{F}_n\right] \le \E\left[\left| X \right| \mid \mathscr{F}_n\right] 
	.\] 
\end{proof}

\begin{ex}
	Let $\{X_j\}_{j \ge 1}$ be independent random variables with mean $\mu$. Let $m < n$. Then
	\begin{align*}
		\E\left[X_1+\ldots+X_n \mid \mathscr{F}_m\right] &= \E\left[X_1 +\ldots+X_m \mid \mathscr{F}_m\right] + \E\left[X_{m+!} + \ldots.+X_n \mid \mathscr{F}_n \right]  \\
		&= X_1+\ldots+X_{m+1} + \E\left[X_1 + \ldots+X_n\right] \\
		&= X_1+\ldots+X_n + \left( n-m \right)\mu
	.\end{align*}
\end{ex}

\begin{ex}
	Let $\{X_j\}_{j\ge 1}$ be i.i.d. random variables with mean $0$ and variance $\sigma^2 > 0$. Suppose $m < n$ and let $S_n = X_1 + \ldots+X_n$ for notational clarity. Then
	\begin{align*}
		S_n^2 &= \left( S_m + S_n - S_m \right)^2 \\
			  &= S_m^2 + (S_n-S_m)^2 + 2S_m(S_n - S_m)
	.\end{align*}
	Then we can write
	\begin{align*}
		\E\left[\left( X_1+\ldots+X_n \right)^2 \mid \mathscr{F}_m\right] &= \E\left[S_m^2 \mid \mathscr{F}_m\right] + \E\left[\left( S_n - S_m \right)^2 \mid \mathscr{F}_m\right] + 2 \E\left[S_m(S_n-S_m) \mid \mathscr{F}_m\right]  \\
		&= S_m^2 + \E\left[(S_n-S_m)^2\right] + 2S_m \E\left[S_n - S_m\right] \\
		&= S_m^2 + (n-m)\sigma^2
	,\end{align*} which is a function of $X_1,\ldots,X_m$, which aligns with what we want from the expectation conditional on $\mathscr{F}_m$.
\end{ex}

\subsection{Introduction to Martingales}

Let $\{Y_j\}_{j\ge_0}$ be random variables and $\mathscr{F}_n$ be the information contained in $Y_0,\ldots,Y_n$. We call $\mathscr{F}_n$ a \emph{filtration}.

Recall that for a random variable $X$ taking values in $\R$ with finite expectation, $\E\left[X \mid \mathscr{F}_n\right] $ is the unique random variable in $\R$ such that 
\begin{enumerate}[(i)]
	\item it is $\mathscr{F}_n$-measurable (a function of $Y_0,\ldots,Y_n$)
	\item for every random variable $Z$ in $\R$ that is $\mathscr{F}_n$-measurable, \[
		\E\left[Z \E\left[X \mid \mathscr{F}_n\right] \right] = \E\left[ZX\right] 
	.\] 
\end{enumerate}

\begin{defn}[Martingale]
	A sequence $\{M_n\}$ of random variables taking values in $\R$ is a martingale with respect to $\{\mathscr{F}_n\}$ if\marginnote{If we don't specify $\{\mathscr{F}_n\}$, assume $\mathscr{F}_n$ is the information in $M_0,\ldots,M_n$.} 
	\begin{enumerate}[(i)]
		\item $M_n$ is $\mathscr{F}_n$-measurable for every $n$
		\item $\E\left[\left| M_n \right|\right] < \infty$ for every $n$.
		\item For every $m \le n$, \[
			\E\left[M_n \mid \mathscr{F}_m \right] = M_m.\] Alternatively, \[ \E\left[M_n - M_m \mid \mathscr{F}_m\right] = 0
		.\] 
	\end{enumerate}
\end{defn}

\begin{rmk}
	Note that $\E\left[M_m\right] = \E\left[M_0\right]$ for every $n$.
\end{rmk}

\begin{prop}
	Property (iii) in the definition of a martingale is equivalent to \[
		\E\left[M_n \mid \mathscr{F}_{n-1}\right] = M_{n-1}, \quad \forall  n \ge 1
	.\] 
\end{prop}
\begin{proof}
	We see that (iii) implies the given condition by picking $m = n-1$. Assume the given condition. Then
	\begin{align*}
		\E\left[M_n \mid \mathscr{F}_{n-2}\right] &= \E\left[ \E\left[M_n \mid \mathscr{F}_{n-1}\right] \mid \mathscr{F}_{n-2}\right] \quad \text{by the tower property} \\
		&= \E\left[M_{n-1} \mid \mathscr{F}_{n-2}\right] \quad \text{by hypothesis} \\
		&= M_{n-2} \quad \text{by hypothesis again}
	.\end{align*}
	We can iterate this same argument $n-m$ times until we get \[
		\E\left[M_n \mid \mathscr{F}_n\right] = M_m
	.\] 
\end{proof}

\begin{ex}
	Let $\{X_j\}$ be independent random variables in $\R$ with mean $\mu$. Let $\mathscr{F}_n$ be the information in $X_1,\ldots,X_n$ and let 
	\begin{align*}
		S_n &= X_1+\ldots+X_n \\
		M_n &= S_n - \mu n
	.\end{align*}

	We claim that $\{M_n\}$ is a martinagle with respect to $\{\mathscr{F}_n\}$.

	\begin{enumerate}[(i)]
		\item $M_n$ is a function of $X_1,\ldots,X_n$, so it is $\mathscr{F}_n$-measurable. 
		\item Check
			\begin{align*}
				\E\left[\left| M_n \right|\right] &\le \sum\limits_{j=1}^{n} \E\left[\left| X_j \right|\right] + \mu n \\
												  &< \infty
			.\end{align*}
		\item Compute
			\begin{align*}
				\E\left[M_n \mid \mathscr{F}_{n-1}\right] &= \E\left[S_{n-1} + X_n - \mu n \mid \mathscr{F}_{n-1}\right]  \\
				&= \E\left[S_{n-1} \mid \mathscr{F}_{n-1}\right] + \E\left[X_n \mid \mathscr{F}_{n-1}\right] - \mu n \\
				&= S_{n-1} + \mu - \mu n \\
				&= S_{n-1} - \mu(n-1) = M_{n-1}
			.\end{align*}
	\end{enumerate}
\end{ex}

\begin{ex}
	Let $\{X_j\}$ be independent random variables in $\R$ with mean $0$ and variance $\sigma^2$. Let $\{\mathscr{F}_n\}$ and $\{S_n\}$ be defined as in the previous example. Define \[
		M_n = S_n^2 - \sigma^2 n
	.\] We claim $\{M_n\}$ is a martinagle with respect to $\{\mathscr{F}_n\}$.

	\begin{enumerate}[(i)]
		\item $M_n$ is $\mathscr{F}_n$-measurable by definition.
		\item Again by the triangle inequality, 
			\begin{align*}
				\E\left[\left| M_n \right|\right] &< \E\left[S_n^2\right] + \sigma^2 n \\
												  &= \E\left[\sum\limits_{i=1}^{n} \sum\limits_{j=1}^{n} X_i X_j \right] + \sigma^2 n \\
												  &= \sum\limits_{i=1}^{n} \sum\limits_{j=1}^{n} \E\left[X_i X_j\right] + \sigma^2 n  \\
												  &= 2\sigma^2 n			
			,\end{align*}
			where the last equality follows from splitting into cases where $i=j$ and $i\ne j$.
		\item Check that
			\begin{align*}
				\E\left[M_n \mid \mathscr{F}_{n-1}\right] &= \E\left[S_n^2 \mid \mathscr{F}_{n-1}\right] -\sigma^2 n \\
				&= S_{n-1}^2 + \sigma^2 - \sigma^2 n \quad \text{by a calculation from last Thursday} \\
				&= M_{n-1}
			.\end{align*}
	\end{enumerate}
\end{ex}

\begin{ex}
	Let $\{\mathscr{F}_n\}$ be any filtration. Let $X$ be any random variable in $\R$ with $\E\left[\left| X \right|\right] < \infty$. Let \[
		M_n = \E\left[X \mid \mathscr{F}_n\right] 
	.\] We claim $\{M_n\}$ is a martingale.

	\begin{enumerate}[(i)]
		\item $M_n$ is $\mathscr{F}_n$-measurable by definition
		\item Note that
			\begin{align*}
				\E\left[\left| M_n \right|\right] &\le \E\left[\E\left[\left| X \right|\mid \mathscr{F}_n\right] \right] \\
				&= \E\left[\left| X \right|\right] < \infty
			.\end{align*}
		\item Check 
			\begin{align*}
				\E\left[M_n \mid \mathscr{F}_{n-1}\right] &= \E\left[\E\left[X \mid \mathscr{F}_n\right] \mid \mathscr{F}_{n-1}\right]  \\
														  &= \E\left[X \mid \mathscr{F}_{n-1}\right] \quad \text{by the tower property} \\
														  &= M_{n-1}
			.\end{align*}
	\end{enumerate}
\end{ex}

\subsection{Optional Stopping Theorem}

Let $\{M_n\}$ be a martingale with respect to $\{\mathscr{F}_n\}$. Recall that \[
	\E\left[M_n\right] = \E\left[M_0\right] \quad \forall n
.\] What happens if we replace $n$ with a random time?\marginnote{Recall that a random variable $\tau \in \N_0 \cup \{\infty\}\}$ is a stopping time for $\{\mathscr{F}_n\}$ if for every $n \in \N_0$, $\{\tau = n\}$ is $\mathscr{F}_n$-measurable. 
In this definition, we can replace the event $\{\tau = n\}$ with $\{\tau \le n\}$ or $\{\tau > n\}$.}

\begin{ex}
	Is $\E\left[M_\tau\right] = \E\left[M_0\right] $ for a stopping time $\tau$?

	Consider $\{X_n\}$ as the random walk on $\Z$ that starts at $0$, which is a martingale with respect to its own filtration, where $\mathscr{F}_n$ is the information in $X_0,\ldots,X_n$. A random walk is a sum of i.i.d. random variables with mean 0, so as we say before, this is a martingale.

	Define \[
		\tau = \min \{n \ge 0 : X_n = 1\}
	.\] We know $\P\left[\tau < \infty\right] = 1$ because the random walk in the integers is a recurrent Markov chain. By the definition of $\tau$, \[
		\E\left[X_\tau\right] = 1 \ne 0 = \E\left[X_0\right] 
	.\] 
\end{ex}

\begin{thm}
	(Optional Stopping Theorem, version 1) Let $\{M_n\}$ be a martingale with respect to $\{\mathscr{F}_n\}$. Let $\tau$ be a stopping time that is bounded; i.e., \[
		\exists \ \text{non-random } K \in \N_0 \text{ s.t. } \P\left[\tau \le K\right]  = 1
	.\] Then \[
		\E\left[M_\tau\right] = \E\left[M_0\right] 
	.\] 
\end{thm}

\begin{proof}
	We can write \[
		M_\tau = \sum\limits_{n=0}^{K} M_n \mathds{1}_{\{\tau=n\}}
	.\] Taking the conditional expectation with respect to $\mathscr{F}_{K-1}$,
	\begin{align*}
		\E\left[M_\tau \mid \mathscr{F}_{K-1}\right] &= \sum\limits_{n=0}^{K-1} \E\left[M_n \mathds{1}_{\{\tau = n\}} \mid \mathscr{F}_{K-1}\right] + \E\left[M_K \mathds{1}_{\{\tau = K\}} \mid \mathscr{F}_{K-1}\right]  \\
		&= \sum\limits_{n=0}^{K-1} M_n \mathds{1}_{\{\tau=n\}} + \E\left[M_K \mathds{1}_{\{\tau > K-1\}} \mid \mathscr{F}_{K-1}\right]  \\
		&= \sum\limits_{n=0}^{K-1} M_n \mathds{1}_{\{\tau=n\}} + \E\left[M_K \mid \mathscr{F}_{K-1}\right] \mathds{1}_{\{\tau > K-1\}} \\
		&= \sum\limits_{n=0}^{K-1} M_n \mathds{1}_{\{\tau=n\}} + M_{K-1} \mathds{1}_{\{\tau > K-1\}} \\
		&= \sum\limits_{n=0}^{K-2}  M_n \mathds{1}_{\{\tau = n\}} + M_{K-1} \left( \mathds{1}_{\{\tau = K-1\}} + \mathds{1}_{\{\tau > K-1\}} \right) \\
		&= \sum\limits_{n=0}^{K-2} M_n \mathds{1}_{\{\tau = n\}} + M_{K-1} \mathds{1}_{\{\tau > K-2\}}
	.\end{align*}
	Now by the tower property,
	\begin{align*}
		\E\left[M_\tau \mid \mathscr{F}_{K-2}\right] &= \E\left[\E\left[M_\tau \mid \mathscr{F}_{K-1}\right] \mid \mathscr{F}_{K-2}\right]  \\
													 &= \sum\limits_{n=0}^{K-2} \E\left[M_n \mathds{1}_{\{\tau = n\}} \mid \mathscr{F}_{K-2}\right] + \E\left[M_{K-1}  \mathds{1}_{\{\tau > K_2\}} \mid \mathscr{F}_{K-2} \right] \\
													 &= \sum\limits_{n=0}^{K-2} M_n \mathds{1}_{\{\tau - n\}} + M_{K-2}\mathds{1}_{\{\tau > K-2\}} \\
													 &= \sum\limits_{n=0}^{K-2} M_n \mathds{1}_{\{\tau = n\}} + M_{K-2} \mathds{1}_{\{\tau > K-3\}}
	.\end{align*}
	We can iterate on this process a total of $K$ times to see \[
		\E\left[M_\tau \mid \mathscr{F}_0\right] = M_0 \mathds{1}_{\{\tau > -1\}} = M_0
	.\] 
	Taking uncondition expectations of both sides, we use the tower property one more time to get \[
		\E\left[M_\tau\right] = \E\left[M_0\right] 
	.\] 
\end{proof}

Now we wonder if there is a version of optional stopping theorem that applies when $\tau$ is not bounded. Let's see what happens when our bound goes to infinity. Define \[
	\tau \wedge n = \min \{\tau,n\}
,\] which is a stopping time because the minimum of two stopping times is a stopping time. It is bounded because $\tau \wedge n \le n$, so we can apply the Optional Stopping Theorem at $\tau \wedge n$.

We see
\begin{align*}
	\E\left[M_0\right] &= \E\left[M_{\tau \wedge n}\right] \\
	&= \E\left[M_\tau \mathds{1}_{\{\tau \le n\}}\right] + \E\left[M_n \mathds{1}_{\{\tau > n\}}\right]
.\end{align*}

What conditions on $\tau$ allow the first term to go to $\E\left[M_\tau\right] $ and the second to $0$?

\begin{thm}
	(Dominated Convergence Theorem) Let $\{X_n\}$ be random variables in $\R$ such that \[
		\lim\limits_{n \to \infty} X_n = X
	\] with probability $1$. Assume there exists $Y$ such that \[
		\E\left[\left| Y \right|\right] <\infty, \qquad \left| X_n \right| \le \left| Y \right|, \ \forall n
	.\] Then \[
		\E\left[X_n\right] \to \E\left[X\right] 
	.\] 
\end{thm}

It is not always true that limits commute with expectation. 

\begin{ex}
	Consider $\{X_n\}$ to be the random variables such that \begin{align*}
		\P\left[X_n = X^2\right] &= \frac{1}{n^2}, \\
		\P\left[X_n = 0\right] &= 1-\frac{1}{n^2}
	.\end{align*}
	Then 
	\begin{align*}
		\P\left[\exists \ j\ge n \text{ s.t. } X_j \ne 0\right] &\le \sum\limits_{j=n}^{\infty} \frac{1}{j^2} \to 0 \\
		\implies \lim\limits_{n \to \infty} X_n &=0 \quad \text{with probability 1}
	,\end{align*} but \[
		\E\left[X_n\right] = 1
	.\] 
\end{ex}

Now, assume that $\E\left[\left| M_\tau \right|\right] <\infty$ and $\P\left[\tau < \infty\right] = 1$. Then \[
	M_\tau \mathds{1}_{\{\tau \le n\}} \to M_\tau
\] with probability 1 and \[
	\left| M_\tau \mathds{1}_{\{\tau \le n\}} \right| \le \left| M_\tau \right| 
.\] Since by assumption, $\E\left[M_\tau\right] < \infty$, we can apply the Dominated Convergence Theorem to see that
\begin{align*}
	\E\left[M_\tau \mathds{1}_{\{\tau \le n\}}\right] &\to \E\left[M_\tau\right]
.\end{align*}

We have no nice way to deal with the second term in the sum for $\E\left[M_{\tau \wedge n}\right]$, so we put the desired condition into our theorem statement.

\begin{thm}
	\marginnote{This version of the Optional Stopping Theorem is much more useful than the first one because most stopping times are not bounded.}(Optional Stopping Theorem, version 2) Let $\{M_n\}$ be a martingale with respect to $\{\mathscr{F}_n\}$ and $\tau$ be a stopping time. Assume \[
		\P\left[\tau < \infty\right] = 1, \quad \lim\limits_{n \to \infty} \E\left[M_n \mathds{1}_{\{\tau > n\}}\right] = 0, \quad \E\left[\left| M_\tau \right|\right] < \infty
	.\] Then \[
		\E\left[M_\tau\right] = \E\left[M_0\right] 
	.\] 
\end{thm}

\begin{prop}
	\label{prop:Mn-tail}
	Suppose $\{M_n\}_{n \ge 0}$ is a sequence of random variables and $\tau$ is a
	random time with $\P(\tau < \infty) = 1$. Assume that either
	\begin{enumerate}[(i)]
		\item there exists $C < \infty$ such that $\P\big[\,|M_n| \le C\,\big] = 1$
		      for every $n$, or
		\item there exists $C < \infty$ such that $\E\left[M_n^2\right] \le C$ for
		      every $n$.
	\end{enumerate}
	Then
	\[
		\lim_{n \to \infty} \E\big[ M_n \mathds{1}_{\{\tau > n\}} \big] = 0.
	\]
\end{prop}

\begin{proof}
	First suppose there exists $C$ with $\P\big[\,|M_n| \le C\,\big] = 1$ for all
	$n$. Then
	\[
		\big| \E\big[ M_n \mathds{1}_{\{\tau > n\}} \big] \big|
		\le \E\big[ |M_n| \mathds{1}_{\{\tau > n\}} \big]
		\le C \P(\tau > n) \xrightarrow[n \to \infty]{} 0,
	\]
	since $\P(\tau > n) \to 0$ when $\P(\tau < \infty) = 1$.

	Next suppose there exists $C$ with $\E\left[M_n^2\right] \le C$ for all $n$.
	By Cauchy–Schwarz,
	\[
		\big| \E\big[ M_n \mathds{1}_{\{\tau > n\}} \big] \big|
		\le \E\left[M_n^2\right]^{1/2} \P(\tau > n)^{1/2}
		\le C^{1/2} \P(\tau > n)^{1/2} \xrightarrow[n \to \infty]{} 0,
	\]
	again using $\P(\tau > n) \to 0$. This proves the claim.
\end{proof}

\begin{ex}
	(Gambler's ruin via optional stopping) Let $\{X_n\}_{n \ge 0}$ be the simple
	symmetric random walk on $\Z$ with $X_0 = 0$. Fix $a,b \in \N$ and define the
	hitting time
	\[
		\tau := \min\{ n \ge 0 : X_n = -a \text{ or } X_n = b \}
	.\] 
	We previously computed $\P\left[ X_\tau = b \right]$ by solving a linear
	difference equation.\marginnote{This is the same gambler’s–ruin setup as in the
		Markov–chain section; we are now redoing the calculation with martingales.}
	Here we give a much shorter derivation using the optional stopping theorem.

	Since $\{X_n\}$ is a martingale with respect to its natural filtration and
	$\tau$ is a stopping time, we can try to apply optional stopping to $\{X_n\}$.
	We have
	\[
		\P\left[ \tau < \infty \right] = 1
	\]
	by recurrence of the simple random walk, and $|X_\tau| \le \max\{a,b\}$, so
	$\E\left[ \left| X_\tau \right| \right] < \infty$. Thus the hypotheses of the
	optional stopping theorem (version with bounded stopping time / bounded
	martingale) are satisfied, and
	\[
		\E\left[ X_\tau \right] = \E\left[ X_0 \right] = 0
	.\] 
	Write
	\[
		p := \P\left[ X_\tau = b \right]
	\]
	so that $\P\left[ X_\tau = -a \right] = 1-p$. Then
	\begin{align*}
		0 = \E\left[X_\tau\right]
		&= b \P\left[ X_\tau = b \right] + (-a)\P\left[ X_\tau = -a \right] \\
		&= b p - a (1-p).
	\end{align*}
	Solving for $p$ gives
	\[
		\P\left[ X_\tau = b \right] = \frac{a}{a+b}
	,\]
	as before, but now by a one–line martingale argument.
\end{ex}

\begin{ex}
	(Using $X_n^2 - n$ to compute $\E[\tau]$) Keep the same setup: $\{X_n\}$ is
	simple symmetric random walk with $X_0 = 0$, and
	\[
		\tau := \min\{ n \ge 0 : X_n = -a \text{ or } X_n = b \}
	.\] 
	Define
	\[
		M_n := X_n^2 - n
	.\] 
	One checks (as in the earlier example) that $\{M_n\}$ is a martingale with
	respect to the natural filtration.\marginnote{The martingale property follows
	from $\E\left[X_{n+1}^2 \mid \mathscr{F}_n\right]
		= X_n^2 + 1$ for the simple symmetric random walk.}

	We will use the extended version of the optional stopping theorem to show that
	\[
		\E\left[ \tau \right] = ab
	.\] 

	First, as before, $\P\left[ \tau < \infty \right] = 1$ and $X_\tau \in \{-a,b\}$.
	Consider the stopped chain $Y_n := X_{\tau \wedge n}$. Then $\{Y_n\}$ is a
	Markov chain on $\{-a,\ldots,b\}$ with three communicating classes
	\[
		\{-a\}, \qquad \{b\}, \qquad \{-a+1,\ldots,b-1\}
	.\] 
	The first two classes are recurrent (once we hit $-a$ or $b$ we stay there),
	while the interior class is transient. By the general result on exit times from
	a transient class, there exists $c>0$ such that
	\[
		\P\left[ \tau > n \right] \le e^{-cn}, \qquad \forall n \ge 1
	.\] 
	In particular,
	\[
		\E\left[ \tau \right]
		= \sum_{n \ge 0} \P\left[ \tau > n \right]
		< \infty
	\]
	and
	\[
		\E\left[ \left| M_\tau \right| \right]
		\le \E\left[ X_\tau^2 \right] + \E\left[ \tau \right]
		\le \max\{a^2,b^2\} + \E\left[\tau\right]
		< \infty
	.\] 
	Moreover,
	\begin{align*}
		\E\left[ \left| M_n \right| \mathds{1}_{\{\tau > n\}} \right]
		&\le \E\left[ \left( \max\{a,b\}^2 + n \right) \mathds{1}_{\{\tau > n\}} \right] \\
		&= \left( \max\{a,b\}^2 + n \right) \P\left[ \tau > n \right] \\
		&\le \left( \max\{a,b\}^2 + n \right) e^{-cn} \xrightarrow[n \to \infty]{} 0.
	\end{align*}
	Thus the hypotheses of the extended optional stopping theorem are satisfied, and
	\[
		\E\left[ M_\tau \right] = \E\left[ M_0 \right] = 0
	.\] 

	On the other hand $X_\tau \in \{-a,b\}$, so with $p := \P\left[ X_\tau = b \right]$
	we have
	\[
		\E\left[ X_\tau^2 \right]
		= a^2 (1-p) + b^2 p
	.\] 
	From the previous example, $p = a/(a+b)$, so
	\begin{align*}
		0 = \E\left[ M_\tau \right]
		&= \E\left[ X_\tau^2 - \tau \right] \\
		&= a^2 (1-p) + b^2 p - \E\left[\tau\right] \\
		&= a^2 \frac{b}{a+b} + b^2 \frac{a}{a+b} - \E\left[\tau\right] \\
		&= ab - \E\left[ \tau \right] .
	\end{align*}
	Therefore $\E\left[ \tau \right] = ab$.
\end{ex}

\subsection{Stopped Martingales}

\begin{prop}
	Let $\{M_n\}_{n \ge 0}$ be a martingale with respect to a filtration
	$\{\mathscr{F}_n\}$. Let $\tau$ be a stopping time. Then the stopped process
	$\{M_{n \wedge \tau}\}_{n \ge 0}$ is also a martingale with respect to
	$\{\mathscr{F}_n\}$.
\end{prop}

\begin{proof}
	Fix $n \ge 0$. Since $\tau$ is a stopping time, the events
	$\{\tau > n\}$ and $\{\tau = k\}$ for $0 \le k \le n$ are all
	$\mathscr{F}_n$–measurable. We can write
	\[
		M_{n \wedge \tau}
		= M_n \mathds{1}_{\{\tau > n\}}
		+ \sum_{k=0}^{n} M_k \mathds{1}_{\{\tau = k\}}
	,\] 
	which is an $\mathscr{F}_n$–measurable linear combination of
	$\mathscr{F}_n$–measurable random variables. This gives the first property.

	For integrability,
	\[
		\left| M_{n \wedge \tau} \right|
		\le \left| M_n \right| + \sum_{k=0}^{n} \left| M_k \right|
	\]
	and each $M_k$ is integrable by assumption, so
	$\E\left[ \left| M_{n \wedge \tau} \right| \right] < \infty$.

	Finally, take $m < n$ and compute the conditional expectation. On the event
	$\{\tau \le m\}$ we have $n \wedge \tau = m \wedge \tau = \tau$, hence
	\[
		M_{n \wedge \tau} = M_\tau = M_{m \wedge \tau}
	\]
	and
	\[
		\E\left[ M_{n \wedge \tau} \mid \mathscr{F}_m \right]
		= M_\tau = M_{m \wedge \tau}
	\]
	on $\{\tau \le m\}$.

	On the complementary event $\{\tau > m\}$ we have $m \wedge \tau = m$ and
	$n \wedge \tau$ is a bounded stopping time (it takes values in
	$\{0,1,\ldots,n\}$). Applying the bounded optional stopping theorem to
	$\{M_k\}$ and the stopping time $n \wedge \tau$ gives
	\[
		\E\left[ M_{n \wedge \tau} \mid \mathscr{F}_m \right]
		= M_m
	\]
	on $\{\tau > m\}$. Combining the two cases,
	\[
		\E\left[ M_{n \wedge \tau} \mid \mathscr{F}_m \right]
		= M_\tau \mathds{1}_{\{\tau \le m\}} + M_m \mathds{1}_{\{\tau > m\}}
		= M_{m \wedge \tau}
	.\] 
	Thus $\{M_{n \wedge \tau}\}$ is a martingale.
\end{proof}

\subsection{The Martingale Convergence Theorem}

\begin{thm}
	(Martingale Convergence Theorem) Let $\{M_n\}_{n \ge 0}$ be a martingale with
	respect to $\{\mathscr{F}_n\}$. Assume that there exists a constant $C>0$ such
	that
	\[
		\E\left[ \left| M_n \right| \right] \le C, \qquad \forall n \ge 0
	.\] 
	Then
	\[
		\P\left[ M_\infty := \lim_{n \to \infty} M_n \text{ exists} \right] = 1
	.\] 
\end{thm}

\marginnote{In other words, a uniformly $L^1$–bounded martingale converges almost
 surely.  A proof can be given via Doob’s upcrossing inequality and is omitted
 here.}

The hypothesis $\sup_n \E\left[ |M_n| \right] < \infty$ can be somewhat
uncomfortable to check directly.  The following gives two convenient sufficient
conditions.

\begin{prop}
	Let $\{M_n\}$ be a martingale. Suppose that either
	\begin{enumerate}[(i)]
		\item $M_n$ is bounded below (or above) uniformly in $n$, i.e., there exists
		      $K>0$ such that either $M_n \ge -K$ a.s.\ for all $n$, or $M_n \le K$
		      a.s.\ for all $n$; or
		\item the second moments are uniformly bounded: there exists $C>0$ such that
		      $\E\left[ M_n^2 \right] \le C$ for all $n$.
	\end{enumerate}
	Then $\sup_{n} \E\left[ |M_n| \right] < \infty$, so the martingale convergence
	theorem applies.
\end{prop}

\begin{proof}
	First assume that there exists $K>0$ such that $M_n \ge -K$ a.s.\ for all $n$.
	Then
	\begin{align*}
		\E\left[ |M_n| \right]
		&= \E\left[ M_n \mathds{1}_{\{M_n \ge 0\}} \right]
		  - \E\left[ M_n \mathds{1}_{\{M_n < 0\}} \right] \\
		&\le \E\left[ M_n \mathds{1}_{\{M_n \ge 0\}} \right] + K \\
		&\le \E\left[ M_n + K \right] + K \\
		&= \E\left[ M_0 \right] + 2K,
	\end{align*}
	which is a constant independent of $n$. The case when $M_n$ is uniformly
	bounded above is similar (apply the previous argument to $-M_n$).

	Now assume $\E\left[ M_n^2 \right] \le C$ for all $n$. By Cauchy–Schwarz,
	\[
		\E\left[ |M_n| \right]
		\le \E\left[ M_n^2 \right]^{1/2} \E\left[ 1^2 \right]^{1/2}
		\le C^{1/2}
	.\] 
	In both cases $\sup_n \E\left[ |M_n| \right] < \infty$, so the convergence
	theorem applies.
\end{proof}

\subsection{Examples of Limits and Non-Limits}

\begin{ex}
	(The limit can be random) Let $\{X_n\}$ be a simple symmetric random walk on
	$\Z$ with $X_0 = 0$, and let $\tau$ be the first time it hits $-a$ or $b$,
	where $a,b \in \N$. Then $\{X_{n \wedge \tau}\}$ is a martingale, and
	\[
		\left| X_{n \wedge \tau} \right| \le \max\{a,b\}
	\]
	for all $n$. Hence $\sup_n \E\left[ |X_{n \wedge \tau}| \right] < \infty$, and
	by the martingale convergence theorem there exists
	\[
		X_\infty := \lim_{n \to \infty} X_{n \wedge \tau} \quad \text{a.s.}
	\]
	In fact $X_{n \wedge \tau} \to X_\tau$ a.s., so $X_\infty$ is a non–degenerate
	random variable taking values in $\{-a,b\}$.
\end{ex}

\begin{ex}
	(Expectation need not be preserved in the limit) Let $\{X_n\}$ be the random
	walk on $\Z$ starting at $X_0 = 0$, and define
	\[
		\tau := \min\{ n \ge 0 : X_n = 1 \}
	.\] 
	Then $\{X_{n \wedge \tau}\}$ is a martingale and $X_{n \wedge \tau} \le 1$ for
	all $n$. As above, the martingale convergence theorem shows that
	\[
		\lim_{n \to \infty} X_{n \wedge \tau} = X_\tau = 1 \quad \text{a.s.}
	\] 
	However, $\E\left[ X_n \right] = 0$ for all $n$, so
	\[
		\E\left[ X_{n \wedge \tau} \right] = \E\left[ X_0 \right] = 0, \qquad \forall n,
	\]
	while
	\[
		\E\left[ X_\tau \right] = 1
	.\] 
	Thus it is \emph{not} true in general that
	\[
		\E\left[ \lim_{n \to \infty} M_n \right] = \lim_{n \to \infty} \E\left[ M_n \right]
	\]
	for a martingale $\{M_n\}$, even when the limit exists.\marginnote{To pass
		expectations through the limit one typically needs an extra hypothesis such
		as uniform integrability.}
\end{ex}

\begin{ex}
	(Not every martingale converges) 
	\begin{enumerate}[(i)]
		\item Let $\{Y_k\}_{k \ge 1}$ be iid with
		      $\P\left[ Y_k = 1 \right] = \P\left[ Y_k = -1 \right] = 1/2$ and define
		      \[
			      M_n := \sum_{k=1}^{n} k Y_k
		      \]
		      with the natural filtration $\mathscr{F}_n := \sigma(Y_1,\ldots,Y_n)$.
		      Then $\{M_n\}$ is a martingale: the increment $M_{n+1}-M_n =
			      (n+1)Y_{n+1}$ has mean zero and is independent of $\mathscr{F}_n$.
		      One can show that $\{M_n\}$ does not converge almost surely.\marginnote{Heuristically,
			      the increments $(n+1)Y_{n+1}$ do not even tend to $0$, so the partial
			      sums cannot converge.}
		\item Let $\{X_n\}$ be the simple symmetric random walk on $\Z$. We know from
		      the recurrence result that $\{X_n\}$ visits every state infinitely often
		      almost surely. In particular, $X_n$ does not converge almost surely.
	\end{enumerate}
\end{ex}

\subsection{Applications of the Martingale Convergence Theorem}

Recall that \[
	\sum\limits_{n=1}^{\infty} \frac{1}{n} = \infty, \qquad \sum\limits_{n=1}^{\infty} \frac{\left( -1 \right)^{n}}{n} \text{ converges}
.\] 
\begin{prop}
	Let $\{Y_n\}_{n\ge 1}$ be i.i.d. such that \[\P\left[Y_n = 1\right] = \P\left[Y_n = -1\right] = \frac{1}{2}.\]
	The randomly alternating harmonic series $\sum\limits_{j=1}^{\infty} \frac{Y_j}{j} $ converges with probability $1$.
\end{prop}
\begin{proof}
	Define $M_n = \sum\limits_{j=1}^{n} \frac{Y_j}{j}$ and let $\mathscr{F}_n$ be the information in $Y_1,\ldots,Y_n$. Then $\{M_n\}$ is a martingale with respect to $\{\mathscr{F}_n\}$. 
	We see that
	\begin{align*}
		\E\left[M_n^2\right] &= \sum\limits_{i=1}^{n} \sum\limits_{j=1}^{n} \E\left[\frac{Y_i Y_j}{ij}\right]  \\
		&= \sum\limits_{j=1}^{n} \frac{\E\left[Y_j^2\right] }{j^2} \\
		&= \sum\limits_{j=1}^{n} \frac{1}{j^2} \\
		&\le \sum\limits_{j=1}^{\infty} \frac{1}{j^2} = \frac{\pi^2}{6}
	.\end{align*}

	The expectation of $M_n$ is uniformly bounded, so by the Martingale Convergence Theorem, $\lim\limits_{n \to \infty} M_n$ exists with probability 1. Hence the randomly alternative hamonic series converges with probability $1$.
\end{proof}

\begin{ex}
	(Polya's Urn) Suppose we have an urn with red and blue balls. At time $n=0$, we have one of each. At time $n$, we draw one ball from the urn, return it, and add an additional ball of the same color as the ball we just drew.

	Let $X_n$ be the number of red balls at time $n$. $\{X_n\}$ is a time inhomogeneous Markov process. 
	\begin{align*}
		\P\left[X_n = x+1 \mid X_{n-1} = x\right] &= \frac{x}{n+1} \\
		\P\left[X_n = x \mid X_{n-1} = x\right] &= \frac{n+1-x}{n+1}
	.\end{align*}
	Let $M_n = \frac{X_n}{n+2}$. We claim $\{M_n\}$ is a martingale with respect to $\{\mathscr{F}_n\}$ being the information at time $n$.
	$M_n$ is a function of $X_n$ and is always within $[0,1]$, so we just need to check the third condition in the definition of a martingale.
	\begin{align*}
		\E\left[M_n \mid \mathscr{F}_{n-1}\right] &= \frac{X_{n-1}}{n+1} \cdot  \frac{X_{n-1}+1}{n+2} + \frac{n+1-X_{n-1}}{n+1} \cdot \frac{X_{n-1}}{n+2} \\
												  &= \frac{\left( X_{n-1} \right) \left( X_{n-1} + 1 + n + 1 - X_{n-1} \right)}{(n+1)(n+2)} \\
												  &= \frac{X_{n-1}}{n+1} \\
												  &= M_{n-1}
	.\end{align*}
	Since $M_n \in [0,1]$, the Martingale Convergence Theorem implies that $\lim\limits_{n \to \infty} M_n$ exists with probability 1. So the fraction of red balls in the urn will converge to something.
	With some more work, we can actually show that $M_{\infty} \sim  Unif(0,1)$.
\end{ex}

\begin{ex}
	Let $X$ be a random variable with $\E\left[\left| X \right|\right] <\infty$. Let \[
		M_n = \E\left[X \mid \mathscr{F}_n\right] 
	\] for any filtration $\{\mathscr{F}_n\}$ such that $\{M_n\}$ is a martingale with respect to $\{\mathscr{F}_n\}$. 
	Note that
	\begin{align*}
		\E\left[\left| M_n \right|\right] &\le \E\left[\E\left[\left| X \right|\mid \mathscr{F}_n\right] \right] \\
										  &= \E\left[\left| X \right|\right]  \in \R
	,\end{align*} so the Martingale Convergence Theorem applies and $M_{\infty}$ exists. We can roughly say that \[
		M_\infty = \E\left[X \mid \mathscr{F}_{\infty}\right] 
	.\] 
\end{ex}

Suppose $X$ is some number we want to estimate and $\mathscr{F}_n$ is the information in $Y_1,\ldots,Y_n$, which are independent samples from the population.\marginnote{As a concrete example, suppose we want to estimate the proportion of the population who will vote for candidate X, and record random people's plans in $Y_1,\ldots,Y_n$.} We can think of $\E\left[X \mid \mathscr{F}_n\right] $ as the best guess for $X$ using $n$ samples.

We should have that \[
	X = \E\left[X \mid \mathscr{F}_\infty\right] 
.\] 

\subsection{Proof of Martingale Convergence Theorem}

We need a new definition and the Monotone Convergence Theorem for the proof. 

\begin{defn}[Upcrossing]
	An upcrossing of $[a,b]$ is an interval in $\{m,\ldots,n\}$ such that $M_m \le a, M_n \ge b$, and for each $k \in \{m+1,\ldots.,n-1\}$, \[
		M_k \in (a,b)
	.\]
\end{defn}

\begin{thm}
	(Monotone Convergence Theorem) Let $\{X_n\}$ be nonnegative random variables in $\R$ such that \[
		X_{n-1} \le X_n
	\] for every $n$. Then \[
		\lim\limits_{n \to \infty} \E\left[X_n\right] = \E\left[\lim\limits_{n \to \infty} X_n\right] 
	.\] 
\end{thm}



Recall the theorem statement.

\begin{thm}
	(Martingale Convergence Theorem) Let $\{M_n\}_{n \ge 0}$ be a martingale with
	respect to $\{\mathscr{F}_n\}$. Assume that there exists a constant $C>0$ such
	that
	\[
		\E\left[ \left| M_n \right| \right] \le C, \qquad \forall n \ge 0
	.\] 
	Then
	\[
		\P\left[ M_\infty := \lim_{n \to \infty} M_n \text{ exists} \right] = 1
	.\] 
\end{thm}

\begin{proof}
	If $\{M_n\}$ does not converge, then \[
		\liminf\limits_{n \to \infty} M_n < \limsup\limits_{n \to \infty} M_n
	.\] Consequently, there exist rational $a < b$ and a sequence of times \[
		m_1 < n_1 < m_2 < n_2 <\ldots
	\] such that \[
		M_{m_j} \le a, \qquad M_{n_j} \le b
	.\] In other words, there exist rational numbers $a$ and $b$ such that $\{M_n\}$ is below $a$ infinitely often and above $b$ infinitely often. Since the expectation of $\left| M_n \right|$ is finite, we know that the martingale cannot go to infinity; it oscillates.

	Hence, if $\{M_n\}$ does not converge, there exist rational $a,b$ such that there exist infinitely many upcrossings of $[a,b]$.\marginnote{Each time interval from $m_j$ to $n_j$ provides at least one upcrossing of $[a,b]$.}

	We claim that with probability $1$, there exist only finitely many upcrossings of $[a,b]$. Note that if we can prove this, there are finitely many upcrossings of $[a,b]$ for every rational $a<b$, so $\{M_n\}$ must converge (otherwise, we showed that $\{M_n\}$ not converging implies there are infinitely many upcrossings, a contradiction).

	Let $B_0 = 1$. Inductively define \[
		B_n = \begin{cases}
			1, \quad B_{n-1} = 0 \text{ and } M_{n-1} \le a \\
			0, \quad B_{n-1} = 1 \text{ and } M_{n-1} \ge b \\
			B_{n-1}, \quad \text{otherwise}
		\end{cases}
	.\] Intuitively, we think of $M_n$ representing the price of an asset, $B_n$ as an indicator for if we own this asset, and the inductive rule defining $B_n$ as buying the stock when $M_{n-1} \le a$ and selling when $M_{n-1} \ge b$.

	By its definition, $B_n$ is $\mathscr{F}_{n-1}$-measurable. Define \[
		W_n = \sum\limits_{j=1}^{n} B_j (M_j - M_{j-1})
	,\] which is the total earnings up until time $n$. 
	Note that $W_n$ is $\mathscr{F}_n$-measurable. Furthermore,
	\begin{align*}
		\E\left[\left| W_n \right|\right] &\le \sum\limits_{j=1}^{n} \E\left[\left| M_j \right| + \left| M_{j-1} \right|\right]
	,\end{align*} which is the finite sum of finite terms.
	Lastly, we see that
	\begin{align*}
		\E\left[W_n \mid \mathscr{F}_{n-1}\right] &= \sum\limits_{j=1}^{n} \E\left[B_j(M_j - M_{j-1}) \mid \mathscr{F}_{n-1}\right]  \\
		&= \sum\limits_{j=1}^{n-1} B_j \left( M_j - M_{j-1} \right) + \E\left[B_n\left( M_n - M_{n-1} \right)\mid \mathscr{F}_{n-1}\right]  \\
		&= \sum\limits_{j=1}^{n-1} B_{j}\left( M_j - M_{j-1} \right) + B_n \E\left[M_n - M_{n-1} \mid \mathscr{F}_{n-1}\right] \\
		&= \sum\limits_{j=1}^{n-1} B_{j}\left( M_j - M_{j-1} \right) = W_{n-1}
	.\end{align*}

	We have shown that $\{W_n\}$ is a martingale with respect to $\{\mathscr{F}_n\}$. Let $U_n$ be the number of upcrossings of $[a,b]$ before time $n$.

	Note that if $B_n = 0$, then $W_n \ge (b-a)U_n$.\marginnote{This is because during an upcrossing, we make profit of at least $b-a$. So the total profit is at least $b-a$ times the total number of upcrossings. This fact is useful because we are trying to upper bound the number of upcrossings and know that $\E\left[W_n\right] =0$.}

	Now consider when $B_n = 1$ and let $k$ be the last time before $n$ such that $B_k = 0$. Since we bought a stock at time $k+1$, $M_k \le a$ and we must sell at the end of every upcrossing, $U_k = U_n$. 

	Hence,
	\begin{align*}
		W_n &= W_k + W_n - W_k \\
			&\ge (b-a)U_k + \sum\limits_{j=k+1}^{n} B_j (M_j - M_{j-1}) \\
			&= (b-a) U_k + \sum\limits_{j=k+1}^{n} (M_j - M_{j-1}) \\
			&= (b-a)U_n + M_n - M_k \\
			&\ge (b-a)U_n - \left| M_n \right| - \left| a \right
	.\end{align*}

	We have this bound for when $B_n = 1$, and the bound for when $B_n = 0$ is a worse bound. We can then combine the two to says that for every $n$, \[
		W_n \ge (b-a)U_n - \left| M_n \right| - \left| a \right|
	.\] 

	Because $\{W_n\}$ is a martingale,
	\begin{align*}
		0 &= \E\left[W_0\right] \\
		&= \E\left[W_n\right]  \\
		&\ge (b-a) \E\left[U_n\right] - \E\left[\left| M_n \right|\right] - \left| a \right| \\
		&\ge (b-a) \E\left[U_n\right] - C - \left| a \right| \quad\text{by hypothesis} \\
		\implies \E\left[U_n\right] \le \frac{C + \left| a \right|}{b-a}
	.\end{align*}
	We have a uniform bound on the expectation of $U_n$. $\{U_n\}$ is nondecreasing, so we can use the Monotone Convergence Theorem that \[
		\E\left[\lim\limits_{n \to \infty} U_n\right] \le \frac{C + \left| a \right|}{b-a}
	.\] The total number of upcrossings has finite expectation, so there must be finitely many upcrossings with probability 1. As we saw earlier, this proves that the limit of $M_n$ exists with probability 1.
\end{proof}

\section{Brownian Motion}

Brownian motion is a real-valued process $\{B_t\}_{t\ge 0}$ where $B: [0,\infty) \to \R$ is a random continuous function. We want a property such that $B_0=0$ and we have stationary increments: for any $t > s> 0$, $B_t - B_s$ has the same disribution as $B_{t-s}$. We also want independent increments: for any \[
	s_1 \le t_1 \le s_2 \le t_2 \le \ldots \le s_k \le t_k
,\] $B_{t_j} - B_{s_j}$ are independent for all $j \in \{1,\ldots,k\}$.

\subsection{Introduction to Brownian Motion}

\begin{defn}[Gaussian (normal) distribution]
	A real-valued random variable $X$ has the Gaussian distribution if the density of $X$ is \[
		f(x) = \frac{1}{\sqrt{2\pi} \sigma} e^{-\frac{\left( x-\mu \right)^2}{2\sigma^2}}
	,\] where $\mu$ is the mean of $X$ and $\sigma^2$ is the variance.
\end{defn}

\begin{thm}
	Let $B: [0,\infty)\to \R$ be a random continuous function with independent stationary increments such that $B_0 = 0$. Then there exists $\mu \in \R$ and $\sigma^2 > 0 $ such that for every $t > 0$, $B_t \sim N(\mu t ,\sigma^2 t)$, where $\mu$ and $\sigma^2$ uniquely characterize the distribution of $B$.
\end{thm}
We give an informal proof for intuition.	
\begin{proof}
	Note that \[
		B_t = \sum\limits_{j=1}^{n} \left( B_{\frac{tj}{n}} - B_{\frac{t(j-1)}{n}} \right)
	\] for every $n \in \N$. The summands are independent and identically distributed. Furthermore, the summands are small for large $n$.

	By a version of the Central Limit Theorem\marginnote{The proof involves showing that the summands have finite variance. Furthermore, the use of the CLT is complicated since the distributions may depend on $n$.} and sending $n\to \infty$, we get that $B_t$ is asymptotically normal.
\end{proof}

\begin{defn}[Brownian Motion]
	Brownian motion is the process $\{B_t\}_{t\ge 0}$ such that $B_0 = 0$ and
	\begin{enumerate}[(i)]
		\item it is continuous,
		\item (stationary increments) for every $s < t$, $B_t - B_s \sim  N(0,t-s)$, and
		\item for every $s_1\le t_1\le\ldots\le s_k \le t_k$, $B_{t_j} - B_{s_j}$ are independent.
	\end{enumerate}
\end{defn}

\begin{thm}
	Browmian motion exists.
\end{thm}

\begin{proof}
	Take Grad Analysis III.	
\end{proof}

We will prove the following properties of Brownian motion:
\begin{itemize}
	\item Brownian motion is nowhere differentiable, 
	\item For every $t >0$ and $\epsilon > 0$, there exist $s_1,s_2 \in (t,t+\epsilon]$ such that \[
			B_{s_1} < B_t, \qquad B_{s_2} > B_t
	.\] 
	\item $\limsup\limits_{t \to \infty} B_t = \infty, \quad \liminf\limits_{t \to \infty} B_t = -\infty$.
\end{itemize}

\begin{ex}
	Let $B$ be standard Brownian motion. Compute
	\begin{align*}
		\P\left[B_1 \ge 1, B_3 \ge B_1 + 1\right] &= \P\left[B_1 \ge 1, B_3 - B_1 \ge 1\right]  \\
												  &= \P\left[\underbrace{B_1 \ge 1}_{N(0,1)}\right] \P\left[\underbrace{B_3 - B_1 \ge 1}_{N(0,2)}\right]  \\
		&= \left( \int_{1}^{\infty} \frac{1}{\sqrt{2\pi} } e^{-\frac{x^2}{2}}dx  \right) \left( \int_{1}^{\infty} \frac{1}{\sqrt{4\pi} } e^{-\frac{x^2}{3}}dx  \right)
	.\end{align*}
\end{ex}

\subsection{Basic Properties of Brownian Motion}
\begin{prop}
	(Brownian Scaling) Let $B$ be Brownian motion and let $c >0$.\marginnote{It is important that $c$ is nonrandom.}
	Then $\{c^{-\frac{1}{2}} B_{ct}\}_{t \ge 0}$ is a Brownian motion.\marginnote{In particular, if $B$ is standard Brownian motion, then $c^{-\frac{1}{2}}B_{ct}$ is also standard Brownian motion.}
\end{prop}

\begin{proof}
	Immediately, \[
		c^{-\frac{1}{2}}B_{c \cdot 0} = 0
	,\] and the increments $\{c^{-\frac{1}{2}} B_{ct}\}$ are scaled increments of the original Brownian motion, so the increments are still independent.

	To check for standard increments, compute\marginnote{This scaling property tells us that no matter how far we zoom in to the path of Brownian motion, it will continue to look the same as the original Brownian motion.}
	\begin{align*}
		c^{-\frac{1}{2}} \left( B_{ct} - B_{cs} \right) &\sim  c^{-\frac{1}{2}} N(0, c(t-s)) \\
														&= N(0,t-s)
	.\end{align*}
\end{proof}

\begin{prop}
	For\marginnote{This is a weaker version of the (in fact true) statement that with probability 1, $B$ is nowhere differentiable.} any $t \ge 0$, \[
		\P\left[B \text{ is not differentiable at $t$}\right] = 1
	.\]
\end{prop}

\begin{proof}
	Let $\epsilon > 0$ and consider $e^{-\frac{1}{2}} B\left( B_{t+\epsilon} - B_t \right) \sim  N(0,1)$.
	For $c >0$,
	\begin{align*}
		\P\left[\left| \frac{B_{t+\epsilon} - B_t}{\epsilon} \right| > c\right] &= \P\left[\left| \epsilon^{-\frac{1}{2}}N(0,1) \right| > c\right]  \\
		&= \P\left[\left| N(0,1) \right| > c \epsilon^{\frac{1}{2}}\right]			
	,\end{align*}
	which goes to $1$ as $\epsilon \to 0$.

	Thus, \[
		\P\left[\limsup\limits_{\epsilon \to 0} \left| \frac{B_{t+\epsilon} - B_t}{\epsilon} \right| = \infty\right] = 1
	,\] so $B$ is not differentiable at $t$.
\end{proof}

\begin{rmk}
	If $\{X_n\}_{n\in \N}$ is in $S$, $A \subset S$ and $\P\left[X_n \in A\right] = 0$ for each $n$, then \[
		\P\left[\exists n \text{ s.t. } X_n \in A\right] \le \sum\limits_{n=1}^{\infty} \P\left[X_n \in A\right] =0
	.\] 

	However, if we have uncountably many random variables (like the values of Brownian motion), this argument does not hold.
	In particular, for every $t > 0$, \[
		\P\left[B_t = 0\right] = 0
	.\] However, \[
		\P\left[\exists\ t > 0 \text{ s.t. } B_t = 0\right] =1 
	.\] 
\end{rmk}

\newthought{We saw that a random walk} is one of most important basic examples in studying Markov chains and martingales. Now we will show that the random walk on $\Z$ converges to Brownian motion. 
\begin{defn}[Convegence in distribution]
	Let $k \in \N$ and $\{\left( X_1^{n},\ldots,X_k^{n} \right)\}_{n\in \N}$ be random variables in $\R^k$. Let $\left( X_1,\ldots,X_k \right)$ be a random variable in $\R^k$.
	We say that $\left( X_1^{n},\ldots,X_k^{n} \right)$ converges to $(X_1,\ldots,X_k)$ in distribution if for every $a_1<b_1,a_2<b_2,\ldots,a_k<b_k$ such that\marginnote{\texthl{why did we define it this way?} \\ Note that \[
		\P\left[X_j = a_j \text{ or }X_j = b_j\right] = 0
	\] is automatic if $X_j$ has a continuous distribution.} \[
		\P\left[X_j =a_j \text{ or } X_j = b_j\right] = 0
	\] for every $j \in \{1,\ldots,k\}$, then \[
		\P\left[X_1^{n} \in (a_1,b_1),\ldots,X_k^{n} \in (a_k,b_k)\right] \to \P\left[X_1 \in (a_1,b_1),\ldots,X_k \in (a_k,b_k)\right] 
	.\] \marginnote{Convergence in distribution relies only on the distributions themselves, not their relations (e.g. we don't care if they're correlated or independent).}

	Equivalently, we have that $\left( X_1^{n},\ldots,X_k^{n} \right) \xrightarrow{\mathscr{D}} (X_1,\ldots,X_k)$ if for every $f: \R^k \to \R$ that is bounded and continuous, \[
		\E\left[f(X_1^{n},\ldots,X_k^{n})\right] \to \E\left[f(X_1,\ldots,X_k)\right] 
	.\] 
\end{defn}

\begin{thm}
	Let $\{X_j\}_{j\ge 0}$ be the unbiased random walk on $\Z$ such that $X_0 = 0$. Extend $X: [0,\infty) \to \R$ by linear interpolation: define \[
		X_t = (t+j)X_{j+1} + (j+1-t)X_j
	\] for every $t \in \left[ j,j+1 \right]$.\marginnote{This theorem is analogous to the Central Limit Theorem.}

	Then $\{n^{-\frac{1}{2}} X_{nt}\}_{t\ge 0}$ converges in distribution to Brownian motion in the following sense: for any $t_1\le t_2 \le \ldots \le t_k$,\marginnote{There are stronger versions of this statement that relate to uniform convergence and the rate of convergence.} \[
	\left( n^{-\frac{1}{2}}X_{nt_1},\ldots,n^{-\frac{1}{2}}X_{n t_k} \right) \xrightarrow{\mathscr{D}} \left( B_{t_1},\ldots,B_{t_k} \right)
	.\] 
\end{thm}

\begin{proof}
	Recall that $\{X_j - X_{j-1}\}_{j \ge 1}$ are i.i.d.\ with mean $0$ and variance $1$. By the Central Limit Theorem, \[
		\frac{X_{tn}}{n^{\frac{1}{2}}} =  t ^{\frac{1}{2}} \frac{X_{tn}}{\left( tn \right)^{\frac{1}{2}}} \to t ^{\frac{1}{2}} N(0,1) = N(0,t) \sim  B_t
	.\] Similarly, if $s <t$, then \[
		\frac{X_{tn} - X_{sn}}{n^{\frac{1}{2}}} \xrightarrow{\mathscr{D}} B_t - B_s
	\] and for $t_1 < t_2 < \ldots < t_k$ where $t_0 = 0$, \[
		X_{nt_j} - X_{nt_{j-1}}
	\] for $j=1,\ldots,k$ are independent.

	Since our random walk increments have independent increments and each increment converges in distribution to an increment of Brownian motion, \[
		n^{-\frac{1}{2}} \left( X_{nt_1}, X_{nt_2} - X_{nt_1}, \ldots, X_{nt_k} - X_{nt_{k-1}} \right) \xrightarrow{\mathscr{D}} \left( B_{t_1}, B_{t_2} - B_{t-1},\ldots,B_{t_k} - B_{t_{k-1}} \right)
	.\] 

	We have that \[
		n^{-\frac{1}{2}} \left( X_{nt_1},X_{nt_2}\ldots,x_{nt_k} \right)
	\] is a linear function of the previous vector of increments of $X$, while \[
	\left( B_{t_1},B_{t_2},\ldots,B_{t_{k-1}} \right)
	\] is also a linear function of the vector of Brownian increments. The convergence in distribution is preserved, so \[
		n^{-\frac{1}{2}} \left( X_{nt_1},X_{nt_2}\ldots,x_{nt_k} \right) \xrightarrow{\mathscr{D}} 	\left( B_{t_1},B_{t_2},\ldots,B_{t_{k-1}} \right)
	.\] 
\end{proof}

\subsection{Relating Brownian Motion to Markov Chains and Martingales}

Let $\mathscr{F}_n$ be the information contained in $\{B_s\}_{s \le t}$. Define the conditional expectation\marginnote{$\{B_s\}$ can be thought of as a random variable taking values in the set of continuous functions on $[0,t]$.} \[
	\E\left[X \mid \mathscr{F}_t\right] = \E\left[X \mid \{B_s\}_{s\le t}\right] 
.\]

\begin{prop}
	(Markov Property) For any $t >0$, $\{B_{s+t} - B_t\}_{t \ge 0}$ is a Brownian motion independent from $\mathscr{F}_t$.
\end{prop}

\begin{proof}
	By the independent increments property of Brownian motion, 
		$\{B_{s+t} - B_t\}_{s\ge 0}$ is independent from $\mathscr{F}_t$. 

	We check that $\{B_{s+t} - B_t\}$ is continuous because the original Brownian motion is continuous. Furthermore, for any $s_1 < s_2$, \[
		\left( B_{s_2+t} - B_t \right) - \left( B_{s_1 + t} - B_t \right) = B_{s_2 + t} - B_{s_1 + t} \sim N(0,s_2-s_1)
	,\] so the future process has the Guassian increments property.

	Fix $s_1\le t_1 \le \ldots \le s_k \le t_k$. Then 
	\begin{align*}
		\left( B_{t_j + t} - B_t \right) - \left( B_{s_j + t} - B_t \right) = B_{t_j + t} - B_{s_j + t}
	\end{align*} is independent by the independent increments property of the original Brownian motion.
\end{proof}

\begin{defn}[Continuous-time martingale]
	Suppose we have $\{M_t\}_{t\ge 0}$, where for each $t\ge 0$, $M_t$ is a real-valued random variable.\marginnote{The only new thing in this defintion is that $s,t$ are allowed to be positive real numbers instead of just natural numbers.} We say that $\{M_t\}_{t\ge 0}$ is a continuous-time martingale with respect to $\mathscr{F}_t$ if
	\begin{enumerate}[(i)]
		\item $M_t$ is $\mathscr{F}_n$-measurable for all $t$
		\item $\E\left[\left| M_t \right|\right] <\infty$ for all $t$
		\item For all $s < t$, \[
			\E\left[M_t \mid \mathscr{F}_s\right] = M_s
		.\] 
	\end{enumerate}
\end{defn}

\begin{prop}
	$\{B_t\}$ is a martingale with respect to $\mathscr{F}_t$.
\end{prop}

\begin{proof}
	$B_t$ is $\mathscr{F}_t$-measurable. $\E\left[\left| B_t \right|\right] < \infty$ since $B_t \sim N(0,t)$. Finally, for $s < t$,
	\begin{align*}
		\E\left[B_t \mid \mathscr{F}_s\right] &= \E\left[B_t - B_s + B_s \mid \mathscr{F}_s\right]  \\
											  &= B_s + \E\left[B_t - B_s \mid \mathscr{F}_s \right] \\
											  &= B_s + \E\left[B_t - B_s\right] \\
											  &= B_s
	.\end{align*}
\end{proof}

\begin{defn}[Stopping time (Brownian motion)]
A\marginnote{\texthl{is it meant to be a stopping time for the filtration?}} random time $\tau \in [0,\infty]$ is a stopping time for $\{\mathscr{F}_t\}$ f for all $t \ge 0$, the event $\{\tau \le t\}$ is $\mathscr{F}_t$-measurable.\marginnote{Like the discrete time case, we can equivalently define a stopping time with the event $\{\tau > t\}$.}
\end{defn}

\begin{ex}
	Let $\tau = \min \{t\ge 0 : B_t = a\}$ for $a \in \R$. Then $\tau$ is a stopping time. 

	$\tau' = \max \{t\le 1 : B_t = 0 \}$ is not a stopping time.
\end{ex}
\begin{prop}
	(Optional stopping theorem) Let $\{M_t\}$ be a continuous martingale with respect to $\mathscr{F}_t$. Let $\tau$ be a stopping time and assume that there exists a nonrandom constant $C > 0$ such that\marginnote{We don't assume that the martingale is bounded --- just that it is bounded by $C$ before $\tau$.} \[
		\P\left[\tau < \infty, \left| M_t \right|\le C, \forall\ t \le \tau\right] = 1
	.\] Then \[
		\E\left[M_\tau\right] = \E\left[M_0\right] 
	.\] 
\end{prop}

\begin{proof}
	First assume that there exists $k \in \N$ such that \[
		\P\left[\tau \in 2^{-k} \N_0\right] = 1
	.\] Let $\tilde{M}_n = M_{2^{-k} n}$, which is a martingale with respect to $\{\mathscr{F}_{2^{-k}n}\}_{n \ge 0}$.

	If $n > m$, then by putting $t = 2^{-k}n$ and $s = 2^{-k}m$, we see \[
		\E\left[M_{2^{-k}n} \mid \mathscr{F}_{2^{-k}m}\right] = M_{2^{-k}m}
	.\] Furthermore, $2^{k}\tau$\marginnote{Since $\tau \in 2^{-k}\N_0$, $2^k \tau \in N_0$.} is a stopping time for $\mathscr{F}_{2^{-k}n}$.

	By the optional stopping theorem on $\{\tilde{M}_n\}$, we get \[
		\E\left[\tilde{M}_{2^k \tau}\right] = \E\left[M_\tau\right]  = \E\left[M_0\right] 
	.\] 

	Now for a general stopping time $\tau$, put $\tau_k = 2^{-k} \left\lceil 2^{k}\tau \right\rceil $ as the smallest integer multiple of $2^{-k}$ which is at least $\tau$. Note that $\tau_k$ is a stopping time $2^{-k}\N_0$, and \[
		\E\left[M_{\tau_k}\right] = \E\left[M_0\right] 
	.\] Since $\lim\limits_{k \to \infty} \tau_k \to \tau$, by continuity of $M$, we know $M_{\tau_k} \to M_\tau$. By some omitted technical argument, this allows us to say \[
		\E\left[M_\tau\right] = \E\left[M_0\right] 
	,\] concluding the proof.
\end{proof}

\begin{prop}
	Let $a,b > 0$ and $\tau = \min \{t \ge 0 : B_t = -a \text{ or } B_t = b\}$.\marginnote{$\tau$ is the first time the Brownian motion started from $0$ exits $(-a,b)$.} Then \[
		\P\left[B_\tau = b\right]  = \frac{a}{a+b}
	.\] 
\end{prop}
\begin{proof}
	We first check that $\tau$ is finite with probability 1:
	\begin{align*}
		\P\left[\tau > t\right] &\le \P\left[B_t \in [-a.b]\right] \\
								&= \P\left[t ^{\frac{1}{2}} N(0,1) \in [-a,b]\right]  \\
								&= \P\left[N(0,1) \in \left[ -at ^{-\frac{1}{2}}, b t ^{-\frac{1}{2}} \right]\right]					
	,\end{align*} 
	which goes to $0$ as $t \to \infty$. Hence $\P\left[\tau < \infty\right] = 1$.

	We also have that \[
		\left| B_t \right| \le \max \{a,b\}
	\] for $t \le \tau$. Then we can put $C = \max \{a,b\}$ and apply the optional stopping theorem to get that
	\begin{align*}
		0 = \E\left[B_0\right] &= \E\left[B_\tau\right] \\
		&= -a \P\left[B_\tau = -a\right] + b \P\left[B_\tau = b\right]  \\
		&= -a \left( 1- \P\left[B_\tau = b\right]  \right) + b \P\left[B_\tau = b\right]  \\
		\P\left[B_\tau = b \right] &= \frac{a}{a+b}
	.\end{align*}
\end{proof}

\begin{rmk}
	In this proof, we did not use anything about Brownian motion specifically besides that it's a continuous time martingale and eventually exits the interval $[-a,b]$. Thus, the conclusion holds for any other continuous time martingale with $\P\left[\tau < \infty\right] = 1$.
\end{rmk}

\begin{prop}
	Brownian motion is recurrent; for any $T > 0$ and $x \in R$, \[
		\P\left[\exists\ t > T \text{ s.t. } B_t = x\right] = 1
	.\] Brownian motion returns to its starting position at arbitrarily large times.
\end{prop}

\begin{proof}
	By the Markov property, $\{B_{t + T} - B_T\}_{t\ge 0}$ is Brownian motion and independent from $\mathscr{F}_T$. 
	We want to show that \[
		\P\left[\exists\ t \ge 0 \text{ s.t. } B_{t+T} - B_T = x - B_T\right] = 1
	.\] Consider when $x - B_T < 0$. We apply the hitting time formula to $\{B_{t + T} - B_T\}_{t\ge 0}$ with \[
		a = -\left( x - B_T \right) > 0, \quad b > 0
	\] conditional on $\mathscr{F}_T$. We can replace the constant $a$ with a random variable because the Brownian motion is independent from $\mathscr{F}_T$, which determines $B_T$. Then
	\begin{align*}
		\P\left[\{B_{t + T} - B_T \text{ hits $b$ before } x - B_T \mid \mathscr{F}_T\}\right] = \frac{-\left( x-B_T \right)}{-\left( x-B_t \right)+b}
	,\end{align*} which goes to $0$ as $b \to \infty$.
	
	With high probability with $b$ is large, $\{B_{t  +T} - B_T\}$ hits $\left( x-B_T \right)$ before $b$. In particular, \[
		\P\left[\{B_{t + T} - B_T\} \text{ hits } x - B_T \mid \mathscr{F}_T\right] = 1
.\] By independence of the Brownian motion from $\mathscr{F}_T$, we see \[
		\P\left[\{B_{t + T} - B_T\} \text{ hits } x - B_T\right] = 1
	.\] 
\end{proof}

\begin{prop}
	(Strong Markov property, continuous time) Let $\tau$ be a stopping time. Then \[
		\{B_{s-\tau} - B_\tau\}_{s \ge 0}
	\] is a Brownian motion that is independent from $\mathscr{F}_\tau$.
\end{prop}

\begin{proof}
	(Outline) First assume that there exists $k \in \N$ such that \[
		\P\left[\tau \in 2^{-k}\N_0\right] = 1
	.\] Use a similar argument as in the discrete time case. In general, we can approximate by \[
		\tau_k = 2^{-k} \left\lceil 2^{k}\tau \right\rceil 
	.\] 
\end{proof}

\begin{prop}
	(Reflection principle) For $a > 0$, \[
		\P\left[\max_{0 \le s \le t} B_s \ge a\right] = 2 \P\left[B_t \ge a\right] 
	.\] Consequently, $\max_{0 \le s \le t} B_s$ has the same distribution as $\left| B_t \right|$. \marginnote{Note that \[
		\{\max_{0\le s \le t} B_s\}_{t\ge 0}
	\] generally does \emph{not} have the same distribution as $\{\left| B_t \right|\}_{t\ge 0}.$}
\end{prop}

\begin{proof}
	Let \[
		\tau = \min \{t \ge 0 : B_t = a\}
	.\] Then $\P\left[\tau < \infty\right] =1$. Furthermore, \[
		\{\max_{0\le s \le t} B_s \ge a\} = \{\tau \le t\}
	,\] since if $B_t \ge a$, then $\tau \le t$.

	We see that \[
		\P\left[B_t \ge a\right] = \P\left[\tau \le t\right] \P\left[B_t \ge a \mid \tau \le t\right] 
	.\] Since $\tau \le t$ is $\mathscr{F}_t$-measurable, by the strong Markov property, the conditional distribution of $B_t - B_\tau$ conditional on $\mathscr{F}_\tau$ is $N(0,t-\tau)$. Then, conditional on $\{t \le \tau\}$, \[
		\P\left[B_t - B_\tau > 0 \mid \mathscr{F}_\tau\right] = \frac{1}{2}
	,\] so \[
		\P\left[B_t - \underbrace{B_\tau}_{a} > 0 \mid t > \tau\right] = \frac{1}{2}
	.\] Continuing,
	\begin{align*}
		\P\left[B_t - a > 0 \mid \tau < t\right] &= \frac{1}{2} \\
		\P\left[B_t \ge a\right] &= \frac{1}{2}\P\left[\tau \le t\right] 
	.\end{align*}

	To get the second conclusion of the reflection principle, we see that
	\begin{align*}
		2\P\left[B_t \ge a\right] &= \P\left[\max_{0\le s \le t} B_s \ge a\right]  \\
		\P\left[\left| B_t \right| \ge a\right] &= \P\left[B_t \ge a\right] + \P\left[-B_t \ge a\right]  \\
												&= 2\P\left[B_t \ge a\right] 
	.\end{align*}
\end{proof}

\subsection{Higher-dimensional Brownian Motion}
\begin{defn}[$d$-dimension Gaussian distribution]
	The Gaussian distribution on $\R^{d}$ with mean $0$ and covariance matrix $\sigma^2 I$ is the distribution of $\overline{x} = \left( x_1,\ldots,x_d \right) \in \R^d$, where $x_1,\ldots,x_d$ are i.i.d.\ from $N(0,\sigma^2)$. We write $\overline{x} \sim  N(0,\sigma^2 I)$.

	The density is given by \[
		\prod_{j=1}^{d} \frac{1}{\sqrt{2\pi} \sigma} e^{\frac{-x_j^2}{2\sigma^2}} = \frac{1}{\left( 2\pi \sigma^2 \right)^{\frac{d}{2}}} e^{-\frac{\left| \overline{x} \right|^2}{2\sigma^2}} 
	,\] where \[
		\left| \overline{x} \right| = \sqrt{x_1^2 + \ldots + x_d^2} 
	.\] The density has the property that for any measurable set $A$, \[
		\P\left[\overline{x} \in A\right]  = \int_{A}^{}  \frac{1}{\left( 2\pi \sigma^2 \right)\frac{d}{2}} e^{-\frac{\left| \overline{x} \right|^2}{2\sigma^2}} dx_1\cdots dx_d
	.\] 
\end{defn}

We implicitly used the standard basis vectors of $\R^d$ in this definition. Are $x_1,\ldots,x_d$ still i.i.d.\ with respect to another basis?

\begin{prop}
	(Rotational invariance) Let $\overline{x} \sim N(0, \sigma^2 I)$. Let $Q $ be a $d\times d$ orthogonal matrix. Then \[
		Qx \sim  N(0,\sigma^2 I)
	.\] 
\end{prop}
\begin{proof}
	Since $Q^TQ = I = Q Q^T$, $\det(Q Q^T) = \det I = 1$. But because \[
		\det(Q Q^T) = \det(Q) \det(Q^T) = \det(Q)^2
	,\] we see that $\left| \det(Q) \right| = 1$.

	Orthogonal\marginnote{\texthl{orthonormal?}} matrices preserve norms of vectors, as in \[
		\left| Q\overline{x} \right| = \overline{x}, \quad \forall\ x \in \R^d
	.\] 

	For $A \subset \{ \R^d\}$ measurable,
	\begin{align*}
		\P\left[Q\overline{x} \in A\right] &= \P\left[\overline{x} \in Q^T A\right]  \\
										   &= \int_{Q^T A}^{} \frac{1}{\left( 2\pi\sigma^2 \right)^{\frac{d}{2}} }e^{- \frac{\left| \overline{x} \right|^2}{2\sigma^2}} dx_1\cdots dx_d \\
	.\end{align*}
	Putting $\overline{x} = Q\overline{y}$, we see that 
	\begin{align*}
		dx_1\cdots dx_d &= \left| \det(Q) \right|dy_1\cdots dy_d \\
						&= dy_1\cdots dy_d
	.\end{align*} 

	Substituting, \begin{align*}
		\P\left[Q\overline{X} \in A\right] &= \int_{A}^{} \frac{1}{\left( 2\pi\sigma^2 \right)^{\frac{d}{2}}}e^{- \frac{\overbrace{\left| Q^T \overline{y} \right|^2}^{\left| \overline{y} \right|}}{2\sigma^2}}dy_1\cdots dy_d  \\ 
										   &= \P\left[\overline{X} \in A\right] 
									   .\end{align*}
\end{proof}

\begin{defn}[Standard Brownian Motion]
	Brownian motion on $\R^d$ is the process \[
		\overline{B}_t = \left( B_t^1 , \ldots, B_t^d \right)
	,\] where $B^1,\ldots, B^d$ are i.i.d.\ Brownian motion in $\R$.
\end{defn}

\begin{prop}
	$d$-dimensional Brownian motion is the unique process in $\R^d$ such that $\overline{B}_0 = 0$ and
	\begin{enumerate}[(i)]
		\item is continuous 
		\item $\forall\ s < t$, we have \[
			\overline{B}_t - \overline{B}_s \sim  N\left( 0,(t-s)I \right)
		.\] 
		\item $\forall\ s_1\le t_1 \le \ldots\le s_k \le t_k$, \[
			\overline{B}_{t_j} - \overline{B}_{s_j}
		\] are independent.
	\end{enumerate}
\end{prop}

\begin{proof}
	\newthought{Existence.} $d$-dimensional Brownian motion satisfies properties (i)-(iii).

	\newthought{Uniqueness.} Suppose $\overline{B}$ satisfies these properties. We can use property (ii) to see that for any $s < t$, \[
		B_t^1 - B_s^1, \ldots ,B_t^d - B_s^d
	\] are independent.
	By this and property (iii), we know $B^1,\ldots, B^d$ are independent. Each $B^j$ for $j=1,\ldots,d$ satisfies the defining properties of $1$-dimensional Brownian motion, so $B^1,\ldots,B^d$ are independent $1$-dimensional Brownian motion.
\end{proof}

\begin{prop}
	Let $c>0$. Then $\{c^{-\frac{1}{2}} \overline{B}_{ct}\}_{t\ge 0}$ is $d$-dimensional Brownian motion.
\end{prop}

\begin{prop}
	Let $\mathscr{F}_t$ be the information in $\{\overline{B}_s\}_{s\le t}$ and $\tau$ be a stopping time for $\{\mathscr{F}_t\}$. Then \[
		\{\overline{B}_{s + \tau} - B_\tau\}_{s\ge 0}
	\] is a $d$-dimensional Brownian motion and is independent from $\mathscr{F_\tau}$.
\end{prop}

\begin{proof}
	The proofs for the scaling and strong Markov properties of $d$-dimensional Brownian motion are identical to the $1$-dimensional cases.
\end{proof}

The following property does not have an analog in the $1$-dimensional case.

\begin{prop}
	(Rotational invariance) Let $Q$ be a $d\times d$ orthogonal\marginnote{\texthl{orthonormal?}} matrix. Then $QB$ is a $d$-dimensional Brownian motion.
\end{prop}
\begin{proof}
	$Q\overline{B}$ is continuous because $Q$ is a linear transformation of continuous $\overline{B}$. It has independent increments for the same reason.
	Furthermore,
	\begin{align*}
		Q \overline{B}_t - Q \overline{B}_s &= Q\left( \overline{B}_t - \overline{B}_s \right) \\
											&= Q N\left( 0, (t-s)I \right) \\
											&= N\left(0, (t-s)I\right)
	.\end{align*}
	$Q\overline{B}$ satisfies all the defining properties of $d$-dimensional Brownian motion.
\end{proof}

Proofs for the following properties of $d$-dimensional Brownian motion can be found in \emph{Brownian Motion} by Morters and Peres.

In $\R^2$, we have
\begin{enumerate}[(i)]
	\item self-intersections in every time interval: $\forall\  a < b$, $\exists\ s,t \in \left[ a,b \right]$ such that $\overline{B}_s = \overline{B}_t$ and $s\ne t$.
	\item neighborhood recurrence: $\forall\ \overline{x} \in \R^2$, $\forall\ \epsilon > 0, \forall\ T > 0$, $ \exists\ t > T$ such that $\left| \overline{B}_t - \overline{x} \right| < \epsilon$. Hence the range of $\overline{B}$ is dense in $R^2$. 
	\item For every $x \in \R^2$, \[
		\P\left[\exists\ t > 0 \text{ s.t. } \overline{B}_t = \overline{x}\right] = 0
	.\] 
\end{enumerate}

In $\R^3$, we have
\begin{enumerate}[(i)]
	\item self-intersection in every time interval 
	\item transience: $\lim\limits_{t \to \infty} \left| \overline{B}_t \right| = \infty$\marginnote{While the $3$-dimensional Brownian motion goes to infinity, its projection onto $\R^2$ is neighborhood recurrent! The graph of this projection looks somewhat like a spiral.}
\end{enumerate}

While in $d=2$, there exist points in $\R^2$ that are hit infinitely many times. However, for $d=3$, each point is hit at most twice.\marginnote{\texthl{why?}}

In $\R^4$ and in higher dimensions, we have
\begin{enumerate}[(i)]
	\item no self-intersections, i.e. \[
		\overline{B}_t \ne \overline{B}_s, \quad \forall\ s\ne t
	.\] 
	\item transience
\end{enumerate}

$d$-dimensional random walks also converge to $d$-dimensional Brownian motion in the same way as the $1$-dimensional case.

\subsection{Relating Brownian Motion and the Heat Equation}
\begin{defn}[Laplacian]
	Let $f: \R^d \to \R$ be $C^2$. The Laplacian of $f$ is \[
		\Delta f(\overline{x}) = \sum\limits_{j=1}^{d} \partial_{x_j}^2 f(\overline{x})
	.\] 
\end{defn}

Theorem (\ref{thm:martbm}) motivates the next definition and example.

\begin{defn}[Harmonic]
	We say that $f: \R^d \to \R$ is harmonic if \[
		\Delta f(\overline{x}) = 0, \quad \forall\ \overline{x} \in \R^d
	.\] 
\end{defn}

Harmonic functions solve the heat equation, as we just put $f(\overline{x},t) := f(\overline{x})$.

\begin{ex}
	Let $\overline{B}_t = \left( B_t^1, B_t^2 \right)$ be a $2$-dimensional Brownian motion. Define \[
		f(x,y) = e^{x} \cos{\left(y\right)}
	.\] Then $f$ is harmonic. 
	Let $\tau = \min \{t\ge 0 : \left| \overline{B}_t - (1,0) \right| \ge 5\}$ be the first time the Brownian motion leaves a disk of radius $5$ centered at $(1,0)$.

	Suppose we want to compute \[
		\E\left[e^{B_\tau^1} \cos{\left(B_\tau^2\right)}\right] 
	.\] We cannot do this explicitly. Since $f$ grows at most exponentially, by our theorem, we have that $f(B_t^1, B_t^2)$ is a martingale. 

	$\tau$ is a stopping time and $\left| f(B_t^1,b_t^2) \right| \le e^{5+1}$\marginnote{The $+1$ comes from the disk not being centered at the origin.} for $t \le \tau$. By the optional stopping theorem, \[
		\E\left[f(B_\tau^1)\right] = \E\left[f(0,0)\right] = 1
	.\] 

\end{ex}

\begin{rmk}
	In complex analysis, for $f: \mathbb{C}\to \mathbb{C}$ holomorphic, the real and imaginary parts of $f$ are harmonic functions. This hints at a connection between $2$-dimensional Brownian motion and complex analysis.
\end{rmk}

\begin{thm} \label{thm:matbm}
	Let $f: \R^d \times [0,\infty) \to \R$\marginnote{Think of $f$ as a function of space and time.} be $C^2$. Assume that for all $t > 0$, there exists $c_t > 0$ such that \[
		\left| f(\overline{x},t) \right| \le c_t e^{c_t \left| \overline{x} \right|}
	\] for all $\overline{x} \in \R^d$. Impose the same restriction on growth on all first and second order partial derivatives of $f$ in $\overline{x}$.
	Further suppose that $f$ solves the heat equation: \[
		\partial_t f(\overline{x},t) + \frac{1}{2} \Delta f(\overline{x},t) = 0
	,\] where the Laplacian is taken only in $\overline{x}$ and $t$ is taken as a constant. Let $\overline{B}$ be a $d$-dimensional Brownian motion.

	Then \[
		M_t = f(\overline{B}_t, t)
	\] is a martingale with respect to $\mathscr{F}_t$, the information in $\{\overline{B}_s\}_{s\le t}$.
\end{thm}
\begin{proof}
	We want to verify that $f(\overline{B}_t,t)$ is a martingale.
	Let \[
		\overline{p}(\overline{x}) = \frac{1}{\left( 2\pi t \right)^{\frac{d}{2}}} e^{- \frac{\left| \overline{x} \right|^2}{2t}}
	,\] which is the density of $\overline{B}_t$. Then 
	\begin{align*}
		\E\left[\left| f(\overline{B}_t, t) \right|\right] &= \int_{\R^d}^{} \left| f(\overline{x}, t) \right| p_t(\overline{x}) dx_1\ldots dx_d  \\
														   &< \infty
	,\end{align*}
	which is finite because of our assumptions on $f$.

	Furthermore, $f(\overline{B}_t, t)$ is $\mathscr{F}_t$-measurable since it is a function of $\overline{B}_t$.

	Now we need to check \[
		\E\left[f(\overline{B}_t, t) \mid \mathscr{F}_s\right]  = f(\overline{B}_s, s)
	.\] 
	First, we can verify that \[
		\partial_t p_t(\overline{x}) - \frac{1}{2}\Delta p_t(\overline{x}) = 0
	,\] so the Gaussian density solves the heat equation with the sign flipped.

	We claim that for any $x_0 \in \R^d$ and $t \ge 0$, \[
		\E\left[f(\overline{B}_t + x_0, t)\right] = f(\overline{x}_0,0)
	.\] We can compute
	\begin{align*}
		\partial_t \E\left[f(\overline{B}_t + \overline{x}_0 , t)\right] &= \partial_t \int_{\R^d}^{} f(\overline{x}+\overline{x}_0, t) p_t(\overline{x})dx_1\cdots dx_d  \\
				 &= \partial_t \int_{\R^d}^{} f(\overline{x},t) p_t(\overline{x}-\overline{x}_0) dx_1\cdots dx_d \text{ by change of variables} \\
				 &= \int_{\R^d}^{} \partial_t p_t(\overline{x}-\overline{x}_0) f(\overline{x},t) + p_t(\overline{x},\overline{x}_0) \parital_t f(\overline{x}, t)dx_1\cdots dx_d \\
				 &= \int_{\R^d}^{} \frac{1}{2} \Delta p_t(\overline{x}-x_0) f(\overline{x},t) + p_t(\overline{x}-\overline{x}_0) \parital_t f(\overline{x}, t) dx_1 \cdots dx_d
	.\end{align*}
	We can integrate the first term by parts so the extra terms decay very quickly according to our assumption on $f$.
	\begin{align*}
		\partial_t \E\left[f(\overline{B}_t + \overline{x}_0 , t)\right] &= \int_{\R^d}^{} p_t (\overline{x} - \overline{x}_0) \left( \frac{1}{2}\Delta f(\overline{x},t)  + \parital_t f(\overline{x},t) \right)dx_1\cdots dx_d = 0	
	,\end{align*}
	where the last equality follows from $f$ saitsfying the heat equation.
	By this result, we know that \[
		\E\left[f(\overline{x} + \overline{B}_t, 0)\right] = f(x_0,0)
	.\] By the Markov property, $\{\overline{B}_{t+s} - \overline{B}_s\}$ is a Brownian motion independent from $\mathscr{F}_s$. 

	We can apply this result to the function \[
		f(\overline{x}, t+ S) ,\quad x_0 = \overline{B}_S
	,\] so 
	\begin{align*}
		\E\left[f(\overline{B}_t, t) \mid \mathscr{F}_{s} \right] &= \E\left[f(\overline{B}_t - \overline{B}_s + \overline{B}_s, t-s _ s \mid \mathscr{F}_s\right]  \\
																  &= f(\overline{B}_s,\frac{s})
															  ,\end{align*}
															  so $\{B_t\}$ is a martingale with respect to $\mathscr{F}_t$.

\end{proof}

\begin{ex}
	Let $B$ be a 1-dimensional Brownian motion and $\theta > 0$. Define \[
		f(x,t) = e^{\theta x - \theta^2 \frac{t}{2}}
	,\] so 
	\begin{align*}
		\partial^2_x &= \theta^2 f(x,t) \\
		\partial_t f(x,t) &= \frac{1}{2} \theta^2 f(x,t)
	\end{align*} and $f$ satisfies the heat equation.

	By our theorem, $f(B_t, t)$ is a martingale. One useful direction to go in once we have a martingale is to use the optional stopping theorem. Put \[
		T_a = \min \{t \ge 0: B_t = a\}
	,\] so for all $t \le T_a$, \[
		f(B_t, t) \le \left[ 0,e^{\theta a} \right]
	.\] We can apply the Optional Stopping Theorem (version 1) to find that \[
		\E\left[f(B_{T_a}, T_a)\right] = f(0,0) = 1
	.\] Note that 
	\begin{align*}
		f(B_{T_a}, T_a) = e^{\theta a - \theta^2 \frac{T_a}{2}} \\
		\E\left[e^{-\theta T_\frac{a}{2}}\right] &= e^{-\theta a}
	.\end{align*}

	For $\theta = \sqrt{2 \lambda}$, where $\lambda >0$, \[
		\E\left[e^{-\lambda T_a}\right]= e^{-2\sqrt{2\lambda} a}
	.\] We derived the Laplace transform (moment generating function) of $T_a$.
\end{ex}

\texthl{It is likely that a Laplace transform (in relating Brownian motion to a martingale) will be on the exam.}

\begin{prop}
	Let $B$ be a 1-dimensional Brownian motion. Define \[
		\tau = \max \{t \in [0,1] : B_t = 0\}
,\] which is \emph{not} a stopping time. We have that \[
		\P\left[\tau \le t\right] = \frac{2}{\pi} \arcsin{\left( \sqrt{t}  \right)}
	.\] 
\end{prop}

\begin{proof}
	We will compute $\P\left[\tau > t\right] $.

	Let $b > 0$. Then the probability of $B$ hitting $0$ again after some given some time $t$ at which it has value $b$ is
	\begin{align*}
		\P\left[\tau > t \mid B_t = b\right] &= \P\left[\min_{s \in [0, t-1]} \left( B_{s+t} - B_t \right) < -b \mid B_t = b\right]  \\
											 &= \P\left[\min_{s \in [0,1-t]} B_s < -b\right] \\
											 &= \P\left[\max_{s \in [0,1-t]} B_s > b\right] \\
											 &= 2\P\left[B_{1-t} > b\right] \text{ by reflection principle} \\
											 &= 2 \int_{b}^{\infty} \frac{1}{\sqrt{2\pi (1-t)} } e^{- \frac{x^2}{2(1-t)}}dx \\
	\text{and substituting $\frac{u}{\sqrt{t} } = \frac{x}{\sqrt{1-t} }$} & \text{ so $\left( \frac{1}{\sqrt{t} } \right)du = \left( \frac{1}{\sqrt{1-t} } \right)dx$} \\
																			 &= 2 \int_{\frac{\sqrt{t} b}{\sqrt{1-t} }}^{\infty} \frac{1}{\sqrt{2\pi t}} e^{- \frac{u^2}{2t}} du 
	.\end{align*}

	We can integrate out the conditioning, multiplying the positive part by $2$ to keep $ b>0$:
	\begin{align*}
		\P\left[\tau > t\right] &= \E\left[\P\left[\tau > t \mid B_t\right] \right]  \\
		&= 2 \int_{0}^{\infty} \P\left[\tau > t \mid B_t = b\right] \frac{1}{\sqrt{2\pi t} }e^{-\frac{b^2}{2t}}db  \\
		&= 4 \int_{0}^{\infty} \int_{\frac{\sqrt{t}}{\sqrt{1-t} }b}^{\infty} \frac{1}{2\pi t} e^{- \frac{b^2+ u^2}{2t}}du db
	\end{align*}

	Since we can evaluate a 2-dimensional Gaussian integral through polar coordinates, put $r = \sqrt{b^2+u^2} $ and $\theta = \theta(u,b) \in [0,2\pi)$ with the bounds \[
		\{0 < b< \infty, \sqrt{\frac{t}{1-t}} b < u <\infty\}
	.\] This means that we have the region of integration \[
		\{0 < r< \infty,\arcsin{(\sqrt{t}) } < \theta < \frac{\pi}{2}\}
	.\] 
	\begin{align*}
		\P\left[\tau > t\right] &= \frac{2}{\pi t} \int_{0}^{\infty} \int_{\arcsin{(s\sqrt{t} )}}^{\pi/2} re^{- \frac{r^2}{2t}}d\theta dr  \\
								&= 1-\frac{2}{\pi} \arcsin{\left( \sqrt{t}  \right)}
	.\end{align*}
\end{proof}

\begin{ex}
	Let $b > a > 0$. What is $\P\left[\exists\ t \in [a,b] \text{ s.t. } b_t = 0\right] $?

	Let $\tau = \max \{t \le b : B_t = 0\}$. We want to compute $\P\left[\tau \ge a\right] $.
	We can rescale (and adding in a $b^{\frac{1}{2}}$ term, which is okay because the equality is with $0$) to \[
		\frac{\tau}{b} = \max \{ s \le 1 : b^{\frac{1}{2}} B_{bs} = 0\}
	.\] 
	Then 
	\begin{align*}
		\P\left[\tau > a\right] &= \P\left[\frac{\tau}{b} > \frac{a}{b}\right] \\
								&= 1-\frac{2}{\pi} \arcsin{\sqrt{\frac{a}{b}} }
	.\end{align*}
\end{ex}

\subsection{Open Problems}
Let $\overline{B}$ be a 2-dimensional Brownian motion.
Consider $\overline{B}([0,1])$ and define $U$ to be the connected components of $\R^2 \setminus \overline{B}([0,1])$. We can view $U$ as a graph where $U\sim V$ iff $\partial I \cap \partial V \ne \O$. The open problem asks: is $U$ connected? More specifically, we have the following:
\begin{conj}
	For any $x,y \in \R^2 \setminus \overline{B}\left( [0,1] \right)$, there exists a path from $x$ to $y$ hitting $\overline{B}\left( [0,1] \right)$ only finitely many times. 
\end{conj}
This is not widely believed one way or another. Gwynne thinks it is not true.


Another question is as follows: What is the size of the smallest path in $\overline{B}\left( [0,1] \right)$ from $\overline{B}_0$ to $\overline{B}_1$? It has been proven that $\overline{B}\left( [0,1] \right)$ does not contain any line segment. One way to measure the size is with Hausdorff dimension. Someone showed that there exists a path in $\overline{B}\left( [0,1] \right)$ from $\overline{B}0$ to $\overline{B}1$ with Hausdorff dimension $\frac{5}{4}$.\marginnote{The Brownian motion itself has Hausdorff dimension $2$, so this is a nontrivial upper bound. This path is a Schram-Loewner evolution with $\kappa = 2$.}

The most common belief is the following.

\begin{conj}
	There exists a path of Hausdorff dimension $1$ but there does not exist any path of finite length. In particular, there is no differentiable path connecting the endpoints of the Brownian motion range.
\end{conj}















\end{document}
